\documentclass{amsart}

\usepackage{amssymb,mathtools,array,longtable}
\usepackage[hidelinks,hypertexnames=false]{hyperref}
\hypersetup{pdftitle={The Goldbach exceptional set: an exponent 11/20},pdfauthor={Ruiyi Zhang}}

\makeatletter
\def\LT@makecaption#1#2#3{\LT@mcol\LT@cols c{\hbox to\z@{\hss
\parbox[t]\textwidth{\normalsize\@tempcnta=64 \@makecaption{#2}{#3}}\hss}}}
\makeatother

\newtheorem{theorem}{Theorem}
\newtheorem{lemma}{Lemma}[section]
\newtheorem{proposition}[lemma]{Proposition}
\newtheorem{corollary}[lemma]{Corollary}
\newtheorem{bilinearlemma}{Lemma}
\renewcommand{\thebilinearlemma}{B}
\newtheorem{primeboxlemma}{Lemma}
\renewcommand{\theprimeboxlemma}{C}
\newcounter{fact}
\renewcommand{\thefact}{F\arabic{fact}}
\newenvironment{fact}{\refstepcounter{fact}\par\medskip
\noindent\textbf{(\thefact)}\ }{\par\medskip}
\numberwithin{equation}{section}
\numberwithin{table}{section}
\allowdisplaybreaks[1]

\newcommand{\N}{\mathbb N}
\newcommand{\Z}{\mathbb Z}
\newcommand{\R}{\mathbb R}
\newcommand{\PP}{\mathbb P}
\newcommand{\ind}{\mathbf 1}
\newcommand{\M}{\mathfrak M}
\newcommand{\m}{\mathfrak m}
\newcommand{\Var}{\operatorname{Var}}
\newcommand{\vol}{\operatorname{vol}}
\newcommand{\Si}{\operatorname{Si}}
\newcommand{\conv}{\operatorname{conv}}
\newcommand{\sums}{\sideset{}{^*}\sum}
\newcommand{\ccert}{c_{\mathrm{cert}}}

\begin{document}

\title{The Goldbach exceptional set: an exponent 11/20}
\author{Ruiyi Zhang}
\address{Independent Researcher, Columbia, Maryland, USA}
\email{ryzfsu@gmail.com}
\subjclass[2020]{Primary 11P32; Secondary 11N35, 11M26}
\keywords{Goldbach problem, exceptional set, Harman sieve, Bombieri--Vinogradov theorem, Dirichlet $L$-functions}

\begin{abstract}
We prove $E(X)\ll_\varepsilon X^{11/20+\varepsilon}$ for every
$\varepsilon>0$. Here $E(X)$ counts the positive even integers at most $X$
that are not sums of two primes. A four-level Harman decomposition gives
a signed prime minorant. Exact rational bounds for four discarded families
give a positive lower bound for its logarithmically weighted mass.
We prove a Bombieri--Vinogradov estimate at level $X^{9/20-\eta}$,
with both maxima and a divisor weight. We also bound the exponential sum
of the same weight on the minor arcs. The major arcs use classical
zero-density estimates and the zero-free half-plane $\Re s>7/8$ of
OpenAI \cite[Theorem~1.1]{OA}.
\end{abstract}

\maketitle

\section{Introduction}\label{sec:introduction}

Let $\PP$ denote the set of primes, and put
\begin{equation}\label{eq:exceptional-set}
E(X)=\#\{n\in\N:n\le X,\ 2\mid n,\ n\ne p+p'
\text{ for all }p,p'\in\PP\}.
\end{equation}
The two primes need not be distinct.

\begin{theorem}\label{thm:main}
For every $\varepsilon>0$,
\begin{equation}
E(X)\ll_\varepsilon X^{11/20+\varepsilon}.
\end{equation}
\end{theorem}

Pintz and Alth\"ofer obtain the exponent $7/12$ from the zero-free
theorem (\ref{fact:zero-free}) below, and Pintz's method gives $3/5$ under
$H_{7/8}$. Theorem~\ref{thm:main} improves both exponents:
\[
\frac{11}{20}<\frac7{12}<\frac35.
\]

The hypothesis and the precise citations are stated below.
The $7/8$ zero-free theorem is an external input and is not reproved here.

\subsection{The zero-free hypothesis}\leavevmode\par
\noindent\textbf{Hypothesis $H_{7/8}$.}
For every Dirichlet character $\chi$,
\begin{equation}\label{eq:H}
L(s,\chi)\ne0\qquad(\Re s>7/8).
\tag{$H_{7/8}$}
\end{equation}
This includes the character modulo $1$, whose $L$-function is $\zeta(s)$.
The functional equation then gives
\begin{equation}\label{eq:zero-strip}
\frac18\le\beta\le\frac78
\end{equation}
for every nontrivial zero $\varrho=\beta+i\gamma$.
An imprimitive character has the same nontrivial zeros as its inducing
primitive character.

\begin{theorem}\label{thm:conditional}
Assume $H_{7/8}$. For every $\varepsilon>0$,
\begin{equation}
E(X)\ll_\varepsilon X^{11/20+\varepsilon}.
\end{equation}
\end{theorem}

\begin{fact}\label{fact:zero-free}
\cite[Theorem~1.1]{OA} Every Dirichlet $L$-function, including $\zeta(s)$,
is zero-free in $\Re s>7/8$.
\end{fact}

\begin{fact}\label{fact:comparison}
\cite[Theorem~1 and Corollary~2]{PA} Theorem~1 assumes a common zero-free
half-plane $H_{\theta_0}$. Corollary~2 invokes
\cite[Theorem~1.1]{OA} and gives the exponent $7/12$ without a hypothesis
in its statement. The exponent obtained from Pintz's method under
$H_{7/8}$ is $3/5$; see \cite[Sections~6--8]{Pi}.
\end{fact}

Theorem~\ref{thm:main} follows from Theorem~\ref{thm:conditional} and
(\ref{fact:zero-free}); the short proof is in Section~\ref{sec:proof-main}.

\subsection{The method}

We construct a signed minorant on $I=(X/8,X/4]$.
Every discarded term is nonnegative, and every negative term is retained
or expanded. The loss of mass comes from four discarded families, and
their rational upper bounds give
\begin{equation}\label{eq:intro-mass}
\sum_{m\in I}(\log m)\rho(m)
\ge\left(\frac{\ccert}{8}-o(1)\right)X,
\qquad
\ccert=\frac{342155430899}{2250000000000}>\frac3{20}.
\end{equation}
Only a lower bound is asserted.

Put $\lambda(m)=(\log m)\rho(m)$; this weight is used throughout.
Its distribution estimate keeps both maxima inside the modulus sum
and includes the factor $2^{\omega(d)}$.
The minor-arc estimate is
\[
\sup_{\alpha\in\mathfrak m}|S_\lambda(\alpha)|
\ll X^{31/40+\eta/2}L^{9/2}.
\]
Under $H_{7/8}$ the major arcs contribute at least $0.025X$, and the
resulting dyadic exceptional-set bound is
\[
|\mathcal E_X|\ll\frac{\mathcal S^2L}{X}
\ll X^{31/20+\eta-1}L^{10}
=X^{11/20+\eta}L^{10}=X^{1-t}L^{10}.
\]
The major-arc argument follows \cite[Section~8]{Pi}. We derive the
required mean square from zero-density estimates and do not use
\cite[(8.25)]{Pi}. Heath-Brown's $20/9$ estimate limits this argument to
levels $t<9/20$. The mass, distribution and minor-arc estimates are all
unconditional.

\section{Notation and parameters}\label{sec:notation}

Write $L=\log X$ and $e(x)=e^{2\pi ix}$.
Given $\varepsilon>0$, choose
\begin{equation}\label{eq:parameters}
\begin{gathered}
\eta=\min(\varepsilon/2,1/100),\qquad
T=\frac9{20},\qquad t=\frac9{20}-\eta,\\
\varepsilon_0=\frac\eta{10},\qquad
X_1=X^{1-\varepsilon_0},\qquad
\varepsilon_{\mathrm{HB}}=\frac\eta{90}.
\end{gathered}
\end{equation}
The decomposition level $T$ is fixed, and the analysis level is $t$.
Put
\begin{equation}\label{eq:scales}
z=X^{1/10},\qquad A=X^{9/20},\qquad D=X^{11/20},
\qquad I=(X/8,X/4].
\end{equation}
Then
\begin{equation}\label{eq:scale-identities}
Az=D,\qquad AD=X,\qquad z^4<A<z^5,\qquad z^{10}=X.
\end{equation}
Let $P$ be any integer satisfying
\begin{equation}\label{eq:P-window}
\frac12X^t\le P\le X^t.
\end{equation}
All estimates are uniform in this choice. We keep the three parameters
\begin{equation}\label{eq:Q-parameters}
Q_P=X/P,\qquad Q_B=X^t,\qquad Q_0=\lceil2X/P\rceil
\end{equation}
distinct; they govern arc widths, distribution and Dirichlet approximation.
The choice of $\varepsilon_0$ gives
\begin{equation}\label{eq:density-margin}
1-\varepsilon_0-\frac{20}{9}t=\frac{191}{90}\eta>0.
\end{equation}
For large $X$,
\begin{equation}\label{eq:elementary-ranges}
\begin{gathered}
2P^2<X,\qquad D\le X/P\le2X^{11/20+\eta},\\
P<X^{1/2}<X_1\le X/8,\qquad A^{3/2}=X^{27/40}.
\end{gathered}
\end{equation}
Implied constants may depend on fixed saving exponents and on $\eta$.
They are uniform in all variables within the stated ranges.

Let $P^-(v)$ be the least prime factor of $v$, with $P^-(1)=+\infty$, and put
\[
\psi(v,w)=\ind_{P^-(v)\ge w}.
\]

Thus $\psi(1,w)=1$, and we call $v$ with $\psi(v,w)=1$ \emph{$w$-rough}.
The function $\Omega(v)$ counts prime factors with multiplicity.
In $2^{\omega(d)}$ the function $\omega(d)$ counts distinct prime factors;
elsewhere $\omega(u)$ denotes Buchstab's function:
\begin{equation}\label{eq:buchstab-definition}
\omega(u)=1/u\quad(1\le u\le2),\qquad
(u\omega(u))'=\omega(u-1)\quad(u>2).
\end{equation}
Buchstab's function is continuous on $[1,\infty)$.

For $f:\N\to\mathbb C$, put
\begin{equation}\label{eq:Delta}
\Delta_f(d)=\max_{(a,d)=1}\max_{y\le X}
\left|\sum_{\substack{m\le y\\m\equiv a\pmod d}}f(m)
-\frac1{\varphi(d)}\sum_{\substack{m\le y\\(m,d)=1}}f(m)\right|.
\end{equation}
Both maxima are part of this definition, and subadditivity and cutoff
monotonicity give
\begin{equation}\label{eq:cutoffs}
\Delta_{f\ind_{[1,Y]}}\le\Delta_f,\qquad
\Delta_{f\ind_{(Y_1,Y_2]}}\le2\Delta_f
\end{equation}
for endpoints at most $X$.

Let $\|x\|$ denote the distance to the nearest integer, and define
\begin{equation}\label{eq:arcs}
\begin{aligned}
\M(q,a)&=\{\alpha\in\R/\Z:\|\alpha-a/q\|\le1/(qQ_P)\},\\
\M&=\bigcup_{q\le P}\ \bigcup_{\substack{a\bmod q\\(a,q)=1}}\M(q,a),
\qquad \m=(\R/\Z)\setminus\M.
\end{aligned}
\end{equation}
These arcs are disjoint, since distinct reduced fractions are at distance
at least $1/(qq')$ and $q+q'\le2P<Q_P$.
Once $\lambda$ is defined, put
\begin{equation}\label{eq:exponential-sums}
S(\alpha)=\sum_{X_1<p\le X}(\log p)e(p\alpha),\qquad
S_\lambda(\alpha)=\sum_{m\in I}\lambda(m)e(m\alpha).
\end{equation}
For $\mathfrak B=\M$ or $\m$, write
\[
I_{\mathfrak B}(n)=\int_{\mathfrak B}S(\alpha)S_\lambda(\alpha)e(-n\alpha)\,d\alpha.
\]
Finally,
\[
c_q(v)=\sum_{\substack{a\bmod q\\(a,q)=1}}e(av/q).
\]
A star on a character sum restricts it to primitive characters.

\section{The four-level Harman weight}\label{sec:weight}

\subsection{Two exact identities}

\begin{lemma}\label{lem:elementary-identities}
For every integer $m>z^2$,
\begin{equation}\label{eq:prime-isolation}
\ind_{\PP}(m)=\psi(m,z)
-\sum_{\substack{p\mid m\\p\ge z,\ p^2\le m}}\psi(m/p,p).
\end{equation}
For $k\ge1$ and $z\le y$,
\begin{equation}\label{eq:buchstab-step}
\psi(k,y)=\psi(k,z)
-\sum_{\substack{p\mid k\\z\le p<y}}\psi(k/p,p).
\end{equation}
\end{lemma}

\begin{proof}
If $m$ is prime, the sum in \eqref{eq:prime-isolation} is empty.
If $m$ is composite and $z$-rough, only $p=P^-(m)$ contributes, and this
prime satisfies $p^2\le m$. If $P^-(m)<z$, every term on the right
vanishes; this case includes prime powers.

For \eqref{eq:buchstab-step}, the case $k=1$ is immediate.
Otherwise put $q=P^-(k)$. For $q\ge y$ the sum is empty, and for
$z\le q<y$ only $p=q$ contributes. For $q<z$, every term vanishes.
\end{proof}

Both identities hold when their left sides vanish, so we expand
zero-valued terms as well.

\subsection{The stopping rule}

For $m\in I$, a node of depth $j$ consists of primes
\[
p_1>p_2>\dots>p_j\ge z,\qquad
P_j:=p_1\cdots p_j\mid m,\qquad p_1^2\le m.
\]
It carries $(-1)^j\psi(m/P_j,p_j)$, and expanding it by \eqref{eq:buchstab-step} gives
\begin{equation}\label{eq:node-expansion}
\begin{split}
(-1)^j\psi(m/P_j,p_j)
={}&(-1)^j\psi(m/P_j,z)\\
&+\sum_{\substack{p_{j+1}\mid m/P_j\\z\le p_{j+1}<p_j}}
(-1)^{j+1}\psi(m/P_{j+1},p_{j+1}).
\end{split}
\end{equation}
The first term is a retained rough leaf, denoted $R_j$, and the root is
$R_0(m)=\psi(m,z)$.

At depth $1$, stop when $A\le p_1\le D$, which gives the leaf $B_1$, and
expand when $p_1<A$.
At depth $2$, stop when $A\le P_2\le D$, which gives $B_2$.
Discard the nodes with $P_2>D$, which form $U_2$, and expand every node
with $P_2<A$.
At depth $3$, stop when $A\le P_3\le D$, giving $B_3$, and expand both
ranges $P_3<A$ and $P_3>D$.
At depth $4$, stop when $A\le P_4\le D$, giving $B_4$.
Discard all other nodes; they form $L_4$, $U_4$ and $V_4$ as specified
below.

Every depth-$1$ node has $p_1\le\sqrt{X/4}<D$.
Every depth-$3$ node has $P_2<A$, and therefore
\begin{equation}\label{eq:third-prime}
p_3<p_2<P_2^{1/2}<A^{1/2},\qquad P_3<A^{3/2}=X^{27/40}.
\end{equation}
If $P_3>D$, then every child has $P_4>D$, so $B_4$ and $L_4$ can arise
only from $P_3<A$. There are no nodes of depth $5$.

Table~\ref{tab:families} defines every family. Each entry is summed over
decreasing primes at least $z$ with $P_j\mid m$. The condition
$p_1^2\le m$ is imposed in every nonempty row.
For $p_1<A$, this condition is automatic on $I$:
\begin{equation}\label{eq:C1-automatic}
p_1^2<X^{9/10}<X/8<m
\end{equation}
for $X>8^{10}$.

\begin{table}[htbp]
\caption{Leaves of the decomposition.}\label{tab:families}
\centering\small
\renewcommand{\arraystretch}{1.18}
\begin{tabular}{@{}c p{0.49\textwidth} c c@{}}
Family&Prefix conditions&Residual&Sign\\ \hline
$R_0$&Empty prefix&$\psi(m,z)$&$+$\\
$R_1$&$p_1<A$&$\psi(m/p_1,z)$&$-$\\
$R_2$&$p_1<A$, $P_2<A$&$\psi(m/P_2,z)$&$+$\\
$R_3$&$p_1<A$, $P_2<A$, and $P_3<A$ or $P_3>D$
&$\psi(m/P_3,z)$&$-$\\
$B_1$&$A\le p_1\le D$&$\psi(m/p_1,p_1)$&$-$\\
$B_2$&$p_1<A$, $A\le P_2\le D$&$\psi(m/P_2,p_2)$&$+$\\
$B_3$&$p_1<A$, $P_2<A$, $A\le P_3\le D$
&$\psi(m/P_3,p_3)$&$-$\\
$B_4$&$p_1<A$, $P_2<A$, $P_3<A$, $A\le P_4\le D$
&$\psi(m/P_4,p_4)$&$+$\\ \hline
$U_2$&$p_1<A$, $P_2>D$&$\psi(m/P_2,p_2)$&$+$\\
$L_4$&$p_1<A$, $P_2<A$, $P_3<A$, $P_4<A$
&$\psi(m/P_4,p_4)$&$+$\\
$U_4$&$p_1<A$, $P_2<A$, $P_3<A$, $P_4>D$
&$\psi(m/P_4,p_4)$&$+$\\
$V_4$&$p_1<A$, $P_2<A$, $P_3>D$
&$\psi(m/P_4,p_4)$&$+$\\ \hline
\end{tabular}
\end{table}

\begin{proposition}\label{prop:tree-identity}
For $m\in I$,
\begin{equation}\label{eq:tree-identity}
\begin{split}
\ind_{\PP}(m)={}&R_0-R_1+R_2-R_3-B_1+B_2-B_3+B_4\\
&+U_2+L_4+U_4+V_4.
\end{split}
\end{equation}
\end{proposition}

\begin{proof}
Start with \eqref{eq:prime-isolation} and apply \eqref{eq:node-expansion}
according to the stopping rule. The process terminates at depth $4$, and
Table~\ref{tab:families} lists every terminal term with its sign.
Every step is an exact identity.
\end{proof}

\subsection{The weight and its support}

On $I$, define
\begin{equation}\label{eq:rho-definition}
\begin{split}
\rho&=R_0-R_1+R_2-R_3-B_1+B_2-B_3+B_4\\
&=\ind_{\PP}-U_2-L_4-U_4-V_4.
\end{split}
\end{equation}
Set $\rho=0$ off $I$, and put
\begin{equation}\label{eq:lambda-definition}
\lambda(m)=(\log m)\rho(m),\qquad
F_\lambda=\sum_{m\in I}\lambda(m).
\end{equation}
Every discarded term is nonnegative. Hence
\begin{equation}\label{eq:pointwise-minorant}
\rho(m)\le\ind_{\PP}(m),\qquad
\lambda(m)\le(\log m)\ind_{\PP}(m).
\end{equation}
The weight may be negative. Every nonzero term is $z$-rough, and since
$m\le X/4<z^{10}$, its support satisfies
\begin{equation}\label{eq:Omega-nine}
\Omega(m)\le9.
\end{equation}
A node is determined by at most four distinct prime factors of $m$, and
it contributes at most one retained term. Thus
\begin{equation}\label{eq:weight-sup}
|\rho(m)|\le\sum_{j=0}^4\binom9j=256,
\qquad |\lambda(m)|\le2^8L.
\end{equation}

For $B_j$, write $u=P_j$ and $v=m/u$; the product $u$ determines its
ordered prime prefix. For $B_1$ we have $v>1$, since $p_1\le D<X/8$, and
then $\psi(v,p_1)=1$ forces $v\ge p_1$. Thus $p_1^2\le m$ is automatic
here as well.
Every balanced family has
\begin{equation}\label{eq:balanced-window}
A\le u\le D,\qquad A/8<v\le D/4.
\end{equation}
Its only roughness coupling is $\psi(v,P^-(u))$, and no retained family
is a bare prime sum.

\section{The mass estimate}\label{sec:mass}

\begin{fact}\label{fact:PNT}
\cite[Chapters~7 and~18]{Da} We use the prime number theorem and Mertens'
estimate in the forms
\begin{gather*}
\pi(x)=\frac{x}{\log x}+O\left(\frac{x}{\log^2x}\right),\qquad
\pi(x)=\operatorname{li}(x)+O\left(\frac{x}{\log^3x}\right),\\
\theta(x)=\sum_{p\le x}\log p=x+O(xe^{-c\sqrt{\log x}}),\\
\sum_{y<p\le x}\frac1p
=\log\frac{\log x}{\log y}+O\left(\frac1{\log y}\right)
\qquad(x\ge y\ge2).
\end{gather*}
Here $\operatorname{li}(x)=\int_2^xdt/\log t$.
\end{fact}

\subsection{A Buchstab envelope}

Put $C_\omega=71/125$ and
\begin{equation}\label{eq:omega-envelope}
\bar\omega(u)=
\begin{cases}
1,&u<1,\\
\max(1/u,C_\omega),&u\ge1.
\end{cases}
\end{equation}

\begin{lemma}\label{lem:omega-envelope}
Let $\Omega_0$ be the positive root of $xe^x=1$. Then
\[
\sup_{u\ge2}\omega(u)=\Omega_0<C_\omega.
\]
The function $\bar\omega$ is continuous and non-increasing.
For $u\ge\max(u_0,1)$,
\[
\omega(u)\le\bar\omega(u_0).
\]
Also $0\le\omega\le1$, and $\omega$ is $1$-Lipschitz.
\end{lemma}

\begin{proof}
On $[2,3]$, equation \eqref{eq:buchstab-definition} gives
\[
u\omega(u)=1+\log(u-1)\qquad(2\le u\le3).
\]
With $w=u-1$, differentiate $(1+\log w)/(1+w)$; its maximum occurs at
$\log w=1/w$ and equals $1/w=\Omega_0$.
For $k\ge3$ put $M_k=\max_{[k,k+1]}\omega$. Then, for $k\le u\le k+1$,
\[
u\omega(u)=k\omega(k)+\int_k^u\omega(t-1)\,dt
\le kM_{k-1}+(u-k)M_{k-1}=uM_{k-1}.
\]
Thus $M_k\le M_{k-1}$, and the rational inequality
\begin{equation}\label{eq:omega-rational}
C_\omega\sum_{j=0}^4\frac{C_\omega^j}{j!}
=\frac{733930882351}{732421875000}>1
\end{equation}
proves $\Omega_0<C_\omega$.
The envelope now follows from $\omega(u)=1/u$ on $[1,2]$.
On this interval $|\omega'|\le1$, and for $u>2$,
\[
|\omega'(u)|=\frac{|\omega(u-1)-\omega(u)|}{u}\le\frac12.
\]
Continuity then completes the Lipschitz assertion.
\end{proof}

\subsection{Rough numbers through \texorpdfstring{$u=10$}{u=10}}

Let
\[
\Phi(x,y)=\#\{1\le n\le x:\ p\mid n\Rightarrow p>y\}.
\]
This count includes $1$.

\begin{lemma}\label{lem:rough-asymptotic}
Uniformly for $y\ge2$ and $y\le x\le y^{10}$,
\begin{equation}\label{eq:rough-asymptotic}
\Phi(x,y)=\frac{x\omega(u)-y}{\log y}
+O\left(\frac{x}{\log^2y}\right),
\qquad u=\frac{\log x}{\log y}.
\end{equation}
The implied constant is absolute.
\end{lemma}

\begin{proof}
Let $\mathrm H(k)$ denote \eqref{eq:rough-asymptotic} for $1\le u\le k$.
For $1\le u\le2$ no composite is counted, so
\[
\Phi(x,y)=1+\pi(x)-\pi(y).
\]
Then (\ref{fact:PNT}) and $\omega(u)=1/u$ prove $\mathrm H(2)$.

Assume $\mathrm H(k)$, where $2\le k\le9$, and take $y^k<x\le y^{k+1}$.
Decomposition by the least prime factor gives
\begin{equation}\label{eq:rough-recursion}
\Phi(x,y)=\Phi(x,\sqrt x)
+\sum_{y<p\le\sqrt x}\Phi(x/p,p)+E_1,
\qquad 0\le E_1\le\frac{2x}{y}.
\end{equation}
The error counts the residuals divisible by the selected prime, and it is
$O(x/\log^2y)$. The base case gives
\[
\Phi(x,\sqrt{x})=\frac{x}{\log x}
+O\left(\frac{x}{\log^2y}\right).
\]
For $y<p\le\sqrt x$ the new parameter is at most $k$. Hence
\begin{equation}\label{eq:rough-induction-term}
\Phi(x/p,p)=g(p)-\frac{p}{\log p}
+O\left(\frac{C_kx}{p\log^2p}\right),
\quad
g(s)=\frac{x}{s\log s}\omega\left(\frac{\log x}{\log s}-1\right).
\end{equation}
Here $C_k$ is the induction constant, and by (\ref{fact:PNT}),
\[
\sum_{p\le\sqrt{x}}\frac{p}{\log p}\ll\frac{x}{\log^2x},
\qquad
\sum_{y<p\le\sqrt{x}}\frac{C_kx}{p\log^2p}
\le C_k\frac{x}{\log^2y}
\left(\log 10+O\left(\frac1{\log y}\right)\right).
\]
The function $g$ is absolutely continuous, and
Lemma~\ref{lem:omega-envelope} gives, almost everywhere,
\[
\begin{split}
|g'(s)|
&\le\frac{x}{s^2\log s}\left(1+\frac1{\log s}\right)
+\frac{x}{s\log s}\cdot\frac{\log x}{s\log^2s}\\
&\le\frac{x}{s^2\log s}
\left(1+\frac{1+u}{\log2}\right)
\le\frac{17x}{s^2\log s}.
\end{split}
\]
Stieltjes integration against $\pi=\operatorname{li}+R$ yields
\[
\sum_{y<p\le\sqrt{x}}g(p)
=\int_y^{\sqrt{x}}\frac{g(s)}{\log s}\,ds
+O\left(\frac{x}{\log^3y}\right).
\]
With $s=x^{1/v}$, the integral equals
\begin{align*}
\frac{x}{\log x}\int_2^u\omega(v-1)\,dv
&=\frac{x}{\log x}\bigl(u\omega(u)-1\bigr)\\
&=\frac{x\omega(u)}{\log y}-\frac{x}{\log x}.
\end{align*}
Insert this in \eqref{eq:rough-recursion}. Since $u>2$, the term
$y/\log y$ is absorbed by the error, and this proves $\mathrm H(k+1)$.
The recurrence $C_{k+1}\ll1+C_k(1+\log10)$ gives an absolute $C_{10}$.
\end{proof}

The subtraction of $y$ in \eqref{eq:rough-asymptotic} is needed near
$x=y$, so we retain it in the interval estimate.

\begin{lemma}\label{lem:rough-interval}
Let $w\ge3$ be an integer and $w'=w-1/2$.
Uniformly for $1\le y_1<y_2\le(w')^{10}$,
\begin{equation}\label{eq:rough-interval}
\begin{split}
N(y_1,y_2;w)
&:=\#\{y_1<v\le y_2:v\ge2,\ P^-(v)\ge w\}\\
&=\frac1{\log w'}\int_{\max(y_1,w')}^{y_2}
\omega\left(\frac{\log t}{\log w'}\right)dt
+O\left(\frac{y_2}{\log^2w}\right).
\end{split}
\end{equation}
The integral is defined to be zero when $y_2\le w'$.
\end{lemma}

\begin{proof}
For integer $v$ the conditions $P^-(v)\ge w$ and $P^-(v)>w'$ coincide,
so $N=\Phi(y_2,w')-\Phi(y_1,w')$. If $y_2\le w'$, this is zero.
If $y_1\ge w'$, the two terms $-w'$ cancel exactly.
If $y_1<w'<y_2$, then $\Phi(y_1,w')=1$, and we use
Lemma~\ref{lem:rough-asymptotic} and $\omega(1)=1$.
In both of the last two cases, with $y_a=\max(y_1,w')$,
\[
N(y_1,y_2;w)
=\frac1{\log w'}
\left[y\,\omega\left(\frac{\log y}{\log w'}\right)\right]_{y_a}^{y_2}
+O\left(\frac{y_2}{\log^2w}\right).
\]
Almost everywhere,
\[
\frac{d}{dy}\left(y\,\omega\left(\frac{\log y}{\log w'}\right)\right)
=\omega(u)+\frac{\omega'(u)}{\log w'},
\qquad u=\frac{\log y}{\log w'}.
\]
Now Lemma~\ref{lem:omega-envelope} proves \eqref{eq:rough-interval}.
\end{proof}

\subsection{Absolute mass and square mass}

\begin{lemma}\label{lem:weight-norms}
We have
\begin{equation}\label{eq:rough-support}
\#\{m\le X:P^-(m)\ge z\}
\le1+\frac{10C_\pi X}{L}\sum_{j=0}^93^j\ll X/L,
\end{equation}
where $\pi(v)\le C_\pi v/\log v$ for $v\ge2$.
Consequently,
\begin{equation}\label{eq:weight-norms}
\sum_{m\in I}|\lambda(m)|\ll X,
\qquad \sum_{m\in I}|\lambda(m)|^2\ll XL.
\end{equation}
\end{lemma}

\begin{proof}
For $1<m\le X$ with $P^-(m)\ge z$, write $m=np$ with $p=P^+(m)$.
Then $\Omega(n)\le9$ and $X/n\ge z$, and the map $m\mapsto(n,p)$ is
injective. The count is therefore at most
\[
1+\sum_{j=0}^9
\sum_{\substack{n:\ \Omega(n)=j,\ P^-(n)\ge z\\ X/n\ge z}}
\pi(X/n)
\le1+\frac{10C_\pi X}{L}
\sum_{j=0}^9\left(\sum_{z\le q\le X}\frac1q\right)^j.
\]
By (\ref{fact:PNT}), the inner sum is $\log10+O(1/L)\le3$.
This proves \eqref{eq:rough-support}, and \eqref{eq:weight-sup} then
gives \eqref{eq:weight-norms}.
\end{proof}

\subsection{The four loss regions}

For $\mathcal F\in\{U_2,L_4,U_4,V_4\}$, let $k=2$ for $U_2$ and $k=4$
otherwise. Let $\mathcal R_{\mathcal F}$ consist of its prime prefixes in
Table~\ref{tab:families}, without $p_1^2\le m$.
For a prefix put
\[
\Pi=P_k,\quad w=p_k,\quad x_i=\frac{\log p_i}{L},\quad
s=\sum_{i=1}^k x_i=\frac{\log\Pi}{L},\quad
e=x_k=\frac{\log w}{L}.
\]
Define
\begin{equation}\label{eq:nonzero-tuples}
\mathcal R_{\mathcal F}^*
=\{(p_1,\ldots,p_k)\in\mathcal R_{\mathcal F}:4\Pi w\le X\}.
\end{equation}
Put $\delta=1/10$ and $s_j=x_1+\cdots+x_j$. The four regions are
\begin{equation}\label{eq:loss-regions}
\begin{aligned}
\mathcal D_{U_2}
&=\{\delta\le b<a<\tfrac9{20},\ a+b>\tfrac{11}{20},\ a+2b\le1\},\\
\mathcal D_{L_4}
&=\{\delta\le e<c<b<a,\ s_4<\tfrac9{20}\},\\
\mathcal D_{U_4}
&=\{\delta\le e<c<b<a,\ s_3<\tfrac9{20},\ s_4>\tfrac{11}{20}\},\\
\mathcal D_{V_4}
&=\{\delta\le e<c<b<a,\ a+b<\tfrac9{20},\\
&\hspace{34mm}s_3>\tfrac{11}{20},\ s_4+e\le1\}.
\end{aligned}
\end{equation}
Here $a,b,c,e$ are the ordered coordinates, and every tuple in
\eqref{eq:nonzero-tuples} lies in its corresponding region. For
\[
u(x):=\frac{1-s}{e},
\]
the ranges are
\begin{equation}\label{eq:buchstab-ranges}
\begin{array}{c|cccc}
\mathcal F&U_2&L_4&U_4&V_4\\ \hline
u(x)&<9/2&44/9<u(x)<6&8/3<u(x)<9/2&<7/2.
\end{array}
\end{equation}
For example, $s_4>4e$ and $e<9/80$ on $\mathcal D_{L_4}$.
On $\mathcal D_{U_4}$ we have $s_4<9/20+e$ and $e<3/20$.
These give the two lower bounds in \eqref{eq:buchstab-ranges}, and the
remaining upper bounds follow from \eqref{eq:loss-regions}.

\begin{lemma}\label{lem:loss-reduction}
For each discarded family,
\begin{equation}\label{eq:loss-reduction}
\sum_{m\in I}(\log m)\mathcal F(m)
\le\frac X8\sum_{\mathcal R_{\mathcal F}^{*}}
\frac{\bar\omega(u_*)}{e\,\Pi}(1+o(1))+O(X/L),
\end{equation}
where
\[
u_*:=\max\left(1,
\frac{\log(X/(8\Pi))}{\log(w-\frac12)}\right).
\]
All errors are uniform over the indicated tuples.
\end{lemma}

\begin{proof}
For $U_2$ and $V_4$ we have $\Pi<A^2$, and for $L_4$ we have $\Pi<A$.
For $U_4$, the inequalities $P_3<A$ and $p_4<P_3^{1/3}$ give
$\Pi<A^{4/3}$.
Thus $\Pi<X^{9/10}$ in every case, and for large $X$ the residual
$v>X/(8\Pi)$ is at least $2$. By \eqref{eq:C1-automatic},
\begin{equation}\label{eq:loss-exact-count}
\sum_{m\in I}\mathcal F(m)
=\sum_{\mathcal R_{\mathcal F}}
N\left(\frac X{8\Pi},\frac X{4\Pi};w\right).
\end{equation}
If $4\Pi w>X$, a counted residual would satisfy
$v\ge w>X/(4\Pi)$.
Its count is therefore exactly zero, and we delete these tuples before
making any estimate.

For an actual tuple,
\begin{equation}\label{eq:finite-buchstab-parameters}
\frac{\log(X/(c\Pi))}{\log(w-1/2)}
=u(x)-\frac{\log c}{eL}+O\left(\frac1{zL}\right)
\qquad(c=4,8).
\end{equation}
By \eqref{eq:buchstab-ranges}, $X/(4\Pi)\le(w-1/2)^{10}$ for large $X$.
Apply Lemma~\ref{lem:rough-interval} to the remaining tuples. The
integration interval has length at most $X/(8\Pi)$, and
Lemma~\ref{lem:omega-envelope} gives
\[
N\left(\frac X{8\Pi},\frac X{4\Pi};w\right)
\le\frac1{\log(w-\frac12)}\cdot\frac X{8\Pi}\,
\bar\omega(u_*)+O\left(\frac X{\Pi L^2}\right).
\]
The summed error is $O(X/L^2)$ by (\ref{fact:PNT}).
Multiply by $\log m\le L$ on the nonnegative losses, and use
$L/\log(w-1/2)=(1+o(1))/e$.
\end{proof}

\subsection{Rectangle majorants}

Fix $n\ge1$ and $k\in\{2,4\}$. Put
\[
\mathfrak h:=\frac{9/20-1/10}{n}=\frac7{20n},\quad
K=\prod_{i=1}^k[x_{i,0},x_{i,1}),\quad
x_{i,0}=\frac1{10}+j_i\mathfrak h,\quad
x_{i,1}=x_{i,0}+\mathfrak h,
\]
where $n>j_1\ge\cdots\ge j_k\ge0$.
Every point of a loss region lies in exactly one such cell. Define
\begin{equation}\label{eq:cell-majorant}
u_0(K)=\frac{1-\sum_i x_{i,1}}{x_{k,1}},\qquad
W(K)=\frac{\bar\omega(u_0(K))\mathfrak h^k}
{x_{1,0}\cdots x_{k-1,0}x_{k,0}^2}.
\end{equation}

\begin{lemma}\label{lem:rectangle}
For $x\in K$ with $s+e\le1$,
\[
\frac{\omega(u(x))}{x_1\cdots x_{k-1}e^2}
\le
\frac{\bar\omega(u_0(K))}
{x_{1,0}\cdots x_{k-1,0}x_{k,0}^2}.
\]
\end{lemma}

\begin{proof}
Write $s'=x_1+\cdots+x_{k-1}$, so $u=(1-s')/e-1$.
First raise $e$ to $x_{k,1}$. The value decreases, since
$1-s'\ge2e>0$ at the starting point. Then raise the other coordinates to
their upper endpoints; each has slope $-1/x_{k,1}$.
Thus $u(x)\ge u_0(K)$, and also $u(x)\ge1$.
Use the non-increasing envelope and the lower endpoints in the denominator.
If $u_0(K)<1$, the envelope equals $1$, so the same bound holds.
\end{proof}

For $U_2$, take $n=256$ and retain cells satisfying
\begin{equation}\label{eq:U2-cell-tests}
a_0<9/20,\qquad a_1+b_1>11/20,\qquad a_0+2b_0\le1.
\end{equation}
For $V_4$, take $n=64$ and retain cells satisfying
\begin{equation}\label{eq:V4-cell-tests}
a_0+b_0<9/20,\qquad a_1+b_1+c_1>11/20,\qquad
a_0+b_0+c_0+2e_0\le1.
\end{equation}
Write these cell sets as $\mathcal K_{U_2}$ and $\mathcal K_{V_4}$.
Each exclusion removes only cells disjoint from the corresponding region.
Define
\begin{equation}\label{eq:cell-sums}
\mathbf U_2=\sum_{K\in\mathcal K_{U_2}}W(K),\qquad
\mathbf V_4=\sum_{K\in\mathcal K_{V_4}}W(K).
\end{equation}
The exact rational arithmetic in Appendix~\ref{app:certificate} gives
\begin{equation}\label{eq:cell-certificates}
\begin{gathered}
|\mathcal K_{U_2}|=11093,\qquad
\mathbf U_2\le\frac{732661466777}{10^{12}},\\
|\mathcal K_{V_4}|=18286,\qquad
\mathbf V_4\le\frac{98012508379}{10^{12}}.
\end{gathered}
\end{equation}

\begin{proposition}\label{prop:rectangle-losses}
We have
\begin{equation}
\sum_{m\in I}(\log m)U_2(m)\le(\mathbf U_2+o(1))\frac X8,
\qquad
\sum_{m\in I}(\log m)V_4(m)\le(\mathbf V_4+o(1))\frac X8.
\end{equation}
\end{proposition}

\begin{proof}
Group the tuples in \eqref{eq:loss-reduction} by cell. By \eqref{eq:finite-buchstab-parameters} and Lemma~\ref{lem:rectangle},
\[
u_*\ge u_0(K)-\frac{10\log8}{L}-O\left(\frac1z\right).
\]
Uniform continuity of $\bar\omega$ gives
$\bar\omega(u_*)\le\bar\omega(u_0(K))+o(1)$, and $1/e\le1/x_{k,0}$.
For each coordinate, (\ref{fact:PNT}) gives
\[
\sum_{X^{x_0}\le p<X^{x_0+\mathfrak h}}\frac1p
=\log\frac{x_0+\mathfrak h}{x_0}+O\left(\frac1L\right)
\le\frac{\mathfrak h}{x_0}+O\left(\frac1L\right).
\]
Enlarge the sum to all prime tuples in the box product; every summand is
nonnegative, including in diagonal cells. Each cell then contributes at
most $W(K)+o(1)$. The number of cells is finite and independent of $X$,
so summing proves the proposition.
\end{proof}

\subsection{The two volume bounds}

\begin{lemma}\label{lem:volume-losses}
We have
\begin{equation}\label{eq:volume-losses}
\begin{aligned}
\sum_{m\in I}(\log m)L_4(m)
&\le\left(\frac{71}{115200}+o(1)\right)\frac X8,\\
\sum_{m\in I}(\log m)U_4(m)
&\le\left(\frac{213}{12800}+o(1)\right)\frac X8.
\end{aligned}
\end{equation}
\end{lemma}

\begin{proof}
By \eqref{eq:buchstab-ranges} and \eqref{eq:finite-buchstab-parameters},
$u_*>2$ for these two families. Hence $\bar\omega(u_*)=C_\omega$ and
$1/e\le10$, and Lemma~\ref{lem:loss-reduction} gives the bound
\[
\sum_{m\in I}(\log m)\mathcal F(m)
\le\frac X8\,C_\omega\,10
\sum_{(p_i)\in\mathcal R_{\mathcal F}^{*}}
\prod_{i=1}^4\frac1{p_i}(1+o(1))+O\left(\frac XL\right).
\]
Cover the fixed region by cubes of side $\sigma$ anchored at $1/10$.
By (\ref{fact:PNT}), each cube contributes at most
$10^4\sigma^4+O(1/L)$, and their total volume tends to
$\vol(\mathcal D_{\mathcal F})$ as $\sigma\to0$.
Letting first $X\to\infty$ and then $\sigma\to0$, we get
\[
\sum_{(p_i)\in\mathcal R_{\mathcal F}^{*}}\prod_i\frac1{p_i}
\le10^4\operatorname{vol}(\mathcal D_{\mathcal F})+o(1).
\]
Use the gap coordinates
\[
e=\delta+y_4,\quad c=e+y_3,\quad b=c+y_2,\quad a=b+y_1.
\]
The Jacobian is $1$, with $y_1,y_2,y_3>0$ and $y_4\ge0$.
For $L_4$ the condition is
\[
y_1+2y_2+3y_3+4y_4<\frac9{20}-4\delta=\frac1{20}=:y_L.
\]
Thus
\begin{equation}\label{eq:L4-volume}
\vol(\mathcal D_{L_4})=\frac{(1/20)^4}{576}.
\end{equation}
For $U_4$, put $Y=y_1+2y_2+3y_3$ and $x_U=3/20$; the two conditions
become
\[
Y+3y_4<\frac3{20}=:x_U,
\qquad
Y+4y_4>\frac{11}{20}-4\delta=\frac3{20}=x_U.
\]
For $0<Y<x_U$ the $y_4$-interval has length $(x_U-Y)/12$, and the
density in $Y$ is $Y^2/12$. Hence
\begin{equation}\label{eq:U4-volume}
\vol(\mathcal D_{U_4})
=\int_0^{x_U}\frac{Y^2(x_U-Y)}{144}\,dY
=\frac{(3/20)^4}{1728}.
\end{equation}
Multiplying these volumes by $C_\omega10^5$ gives
$71/115200$ and $213/12800$, respectively.
\end{proof}

\subsection{The mass lower bound}

\begin{proposition}\label{prop:mass}
Put
\begin{equation}\label{eq:ccert}
\begin{split}
\ccert
&=1-\frac{732661466777}{10^{12}}-\frac{71}{115200}
-\frac{213}{12800}-\frac{98012508379}{10^{12}}\\
&=\frac{342155430899}{2250000000000}>\frac3{20}.
\end{split}
\end{equation}
Then
\begin{equation}\label{eq:mass-bound}
F_\lambda\ge\left(\frac{\ccert}{8}-o(1)\right)X,
\qquad |F_\lambda|\le\sum_{m\in I}|\lambda(m)|\ll X.
\end{equation}
\end{proposition}

\begin{proof}
By (\ref{fact:PNT}), the prime contribution is
\[
\theta(X/4)-\theta(X/8)
=\frac X8+O\left(Xe^{-c\sqrt L}\right).
\]
Subtract the four nonnegative losses in \eqref{eq:rho-definition}, using
Proposition~\ref{prop:rectangle-losses}, \eqref{eq:cell-certificates}
and Lemma~\ref{lem:volume-losses}. This retains the exact weight $\log m$
on primes, and on the losses the upper bound $\log m\le L$ has the
required direction. The absolute bound follows from
Lemma~\ref{lem:weight-norms}.
\end{proof}

\section{Duality}\label{sec:duality}

For even $n$, put
\begin{equation}
r_\lambda(n):=\sum_{X_1<p\le X}(\log p)\,\lambda(n-p)
=\int_0^1S(\alpha)S_\lambda(\alpha)e(-n\alpha)\,d\alpha
=I_\M(n)+I_\m(n).
\end{equation}

Both integrals are real, since the arcs are invariant under
$\alpha\mapsto-\alpha$ and the coefficients are real.
If $n$ is not a sum of two primes, then
\begin{equation}\label{eq:exceptional-sign}
r_\lambda(n)\le0.
\end{equation}
This follows term by term from \eqref{eq:pointwise-minorant}.
For $n\in(X/2,X]$ and $m\in I$ we also have
\begin{equation}\label{eq:prime-location}
n-m\in[X/4,7X/8]\subset(X_1,X].
\end{equation}

\begin{proposition}\label{prop:duality}
Let $\mathcal E_X$ be the exceptional even integers in $(X/2,X]$.
Suppose that $I_{\M}(n)\ge c_1X$ for every even $n$ in this interval,
and that $\sup_{\alpha\in\m}|S_\lambda(\alpha)|\le\mathcal S$.
Then
\[
|\mathcal E_X|\le c_1^{-2}\mathcal S^2L/X.
\]
\end{proposition}

\begin{proof}
Put $Z(\alpha)=\sum_{n\in\mathcal E_X}e(n\alpha)$.
By \eqref{eq:exceptional-sign},
\[
c_1X\le I_{\M}(n)\le-I_{\m}(n)\qquad(n\in\mathcal E_X).
\]
Parseval gives $\|Z\|_2^2=|\mathcal E_X|$, and by (\ref{fact:PNT}),
$\|S\|_2^2\le XL$ for large $X$. Squaring proves the assertion.
\end{proof}

\section{Minor arcs}\label{sec:minor}

This section is unconditional. For intervals $J\subset I$ put
\[
T_\rho(\alpha;J)=\sum_{m\in J}\rho(m)e(m\alpha).
\]
We first bound these sums uniformly in $J$.

\subsection{Rational approximation and bilinear estimates}

For $\alpha\in\m$, Dirichlet's approximation theorem gives coprime $a,s$ with
\begin{equation}\label{eq:minor-approximation}
P<s\le Q_0,\qquad
|\alpha-a/s|\le(sQ_0)^{-1}\le s^{-2}.
\end{equation}
A denominator $s\le P$ would place $\alpha$ on a major arc, since
$(sQ_0)^{-1}\le P/(2sX)$. Thus
\begin{equation}\label{eq:minor-sizes}
X/s<X/P\le2X^{11/20+\eta},\qquad
s\le3X/P\le6X^{11/20+\eta},\qquad D\le X/P.
\end{equation}

\begin{fact}\label{fact:reciprocal}
\cite[Lemma~2.2]{C9} If $(a,s)=1$ and $|\alpha-a/s|\le s^{-2}$, then
\begin{align*}
\sum_{u\le U}\min\left(\frac{2X}{u},\frac1{2\|\alpha u\|}\right)
&\ll (X/s+U+s)\log(2Xs),\\
\sum_{h\le H}\min\left(Y,\frac1{\|\alpha h\|}\right)
&\ll(1+H/s)(Y+s\log s).
\end{align*}
\end{fact}

\begin{lemma}\label{lem:type-I-II}
Let $|a_u|,|b_v|\le1$, and let $J\subset(0,X]$ be an interval.
Then
\begin{equation}\label{eq:type-I}
\left|\sum_{u\le U}a_u\sum_{k:uk\in J}e(\alpha uk)\right|
\ll(X/s+U+s)\log(2Xs).
\end{equation}
If $MN\ll X$, then
\begin{equation}\label{eq:type-II}
\left|\sum_{M<u\le2M}\sum_{v\le N}
a_ub_v\ind_J(uv)e(\alpha uv)\right|
\ll X^{1/2}(X/s+M+N+s)^{1/2}L^{1/2}.
\end{equation}
\end{lemma}

\begin{proof}
The inner sum in \eqref{eq:type-I} is a geometric progression with
modulus at most $\min(2X/u,(2\|\alpha u\|)^{-1})$, so
(\ref{fact:reciprocal}) applies.
For \eqref{eq:type-II}, apply Cauchy--Schwarz in $u$. For fixed $v,v'$,
the conditions $uv,uv'\in J$ restrict $u$ to an interval.
Group the off-diagonal terms by $h=|v-v'|$; the second estimate in
(\ref{fact:reciprocal}) gives
\begin{align*}
|T|^2
&\ll MN\left(M+\sum_{1\le h<N}
\min\left(2M,\frac1{\|\alpha h\|}\right)\right)\\
&\ll MN\left(M+\left(1+\frac Ns\right)(M+s\log s)\right)\\
&\ll MN\left(\frac{MN}s+M+N+s\right)L,
\end{align*}
where $T$ denotes the left side of \eqref{eq:type-II}.
Since $MN\ll X$, this proves \eqref{eq:type-II}.
\end{proof}

By \eqref{eq:minor-sizes}, the estimates needed below are
\begin{align}
U\le X^{27/40}&:\quad \text{\eqref{eq:type-I}}\ll X^{27/40}L,
\label{eq:type-I-size}\\
M,N\ll X/P&:\quad \text{\eqref{eq:type-II}}
\ll X^{31/40+\eta/2}L^{1/2}.
\label{eq:type-II-size}
\end{align}
The first of these uses $11/20+\eta\le0.56<27/40$.

\subsection{Exact separation of comparisons}

Let $K=2\lceil X\rceil+1$, and let $H$ be the indicator of
$\{0,\ldots,\lceil X\rceil\}$ modulo $K$.
Fourier inversion gives, for integers $x\in[-X,X]$,
\begin{equation}\label{eq:Fourier-comparison}
\ind[x\ge0]=\sum_{g\bmod K}\widehat H(g)e(gx/K),
\qquad \sum_g|\widehat H(g)|\le2+\log K\ll L,
\end{equation}
where $\widehat H(g)=K^{-1}\sum_{h=0}^{\lceil X\rceil}e(-gh/K)$.
The norm bound follows from the geometric-sum bound
$|\widehat H(g)|\le\min(1,(2K\|g/K\|)^{-1})$.
Consequently, for $2\le u,v\le X$,
\begin{equation}\label{eq:separate-roughness}
\psi(v,P^-(u))
=\sum_g\widehat H(g)e(gP^-(v)/K)e(-gP^-(u)/K).
\end{equation}
If $F(u)\in[0,X]$ is integer-valued, then
\begin{equation}\label{eq:separate-cutoff}
\ind[w\le F(u)]
=\sum_g\widehat H(g)e(gF(u)/K)e(-gw/K)
\qquad(1\le w\le X).
\end{equation}
These identities include equality at each endpoint, and each separated
comparison costs $O(L)$.

\subsection{The stopped Legendre identity}

\begin{lemma}\label{lem:stopped-Legendre}
Let $r<A$ be a positive integer, and let $k\ge1$. Then
\begin{equation}\label{eq:stopped-Legendre}
\begin{split}
\psi(k,z)
={}&\sum_{\substack{d\mid k,\ d\mid\Pi(z)\\rd<A}}\mu(d)\\
&+\sum_{\substack{d\mid k,\ d\mid\Pi(z),\ d>1\\
rd\ge A>rd/P^-(d)}}
\mu(d)\psi(k/d,P^-(d)),
\end{split}
\end{equation}
where $\Pi(z)=\prod_{p<z}p$.
Every term in the second sum has $A\le rd<Az=D$.
\end{lemma}

\begin{proof}
Iterate
\[
\psi(k',y)=1-\sum_{\substack{p<y\\p\mid k'}}\psi(k'/p,p).
\]
The selected primes decrease strictly, so their product $d$ is
squarefree. Expand a node precisely when $rd<A$.
At a stopped node, $q=P^-(d)<z$ and
$rd/q<A\le rd<Aq<Az$.
The terminal terms give \eqref{eq:stopped-Legendre}.
The rule depends only on $r,d$, and no restriction is placed on repeated
prime factors of the residual.
\end{proof}

\subsection{The eight retained families}

\begin{proposition}\label{prop:minor-rho}
Uniformly for $\alpha\in\m$ and intervals $J\subset I$,
\begin{equation}
|T_\rho(\alpha;J)|\ll X^{31/40+\eta/2}L^{7/2}.
\end{equation}
\end{proposition}

\begin{proof}
We use the eight terms in \eqref{eq:rho-definition}; the four discarded
families require no exponential-sum estimate.

\noindent\textit{The families $R_0,R_1,R_2$.}
For $j=0,1,2$ the sums have the form
\begin{equation*}
\sum_{m\in J}R_j(m)e(\alpha m)
=\sum_r c_r\sum_{k:\ rk\in J}\psi(k,z)e(\alpha rk),
\end{equation*}
where $r=P_j<A$, with $r=1$ when $j=0$.
Here $c_r\in\{0,1\}$ depends only on $r$, which is $z$-rough.
Apply Lemma~\ref{lem:stopped-Legendre} and put $u=rd$. The rough and
smooth parts of $u$ determine $r,d$ uniquely, so the resulting
coefficients have modulus at most $1$. The free terms have $u<A$ and
contribute
$O(X^{11/20+\eta}L)$ by \eqref{eq:type-I}.
The stopped terms are
\[
\sum_{A\le u<D}a'_u\sum_{k':\ uk'\in J}
\psi(k',P^-(u))e(\alpha uk'),
\qquad |a'_u|\le1,\quad P^-(u)=P^-(d).
\]
Here $k'>X/(8D)>2$. Split $u$ into $O(L)$ intervals $(M,2M]$ with
$A/2\le M<D$; then $k'\le X/(4M)$.
By \eqref{eq:separate-roughness} and \eqref{eq:type-II-size}, the total is
$O(X^{31/40+\eta/2}L^{5/2})$.

\noindent\textit{The family $R_3$.}
Write $r=P_2=p_1p_2<A$ and $w=p_3$. The prime conditions are
\[
p_1>p_2>w\ge z,\qquad rw<A\ \text{or}\ rw>D.
\]
Apply Lemma~\ref{lem:stopped-Legendre} with anchor $r=P_2$ to the
residual $\psi(k,z)$ in $m=rwk$; the anchor is not $P_3$.
For free terms the outer variable is $U=rwd$. Its rough part determines
$p_1,p_2,p_3$, and its smooth part determines $d$, so the coefficient is
bounded by $1$. The inner variable is free on $J/U$.
Since $w^2<p_2^2<r<A$,
\[
U=(rd)w<A^{3/2}=X^{27/40}.
\]
These terms contribute $O(X^{27/40}L)$, for both conditions on $rw$.

For stopped terms, put $u=rd\in[A,D)$, $q=P^-(d)<z$, and $V=wk'$.
Since $w\ge z>q$, we have $\psi(k',q)=\psi(V,q)$, and the sum becomes
\begin{equation}\label{eq:R3-stopped}
\sum_{A\le u<D}a_u\sum_{V:uV\in J}
\psi(V,P^-(u))
\sum_{\substack{w\mid V,\ w\in\PP,\ w\ge z}}
\vartheta(u,w)e(\alpha uV).
\end{equation}
Here $|a_u|\le1$ includes the stopping condition, and the rough part of
$u$ is $r=p_1p_2$. Also
\[
\vartheta(u,w)=\ind[w<p_2]
\bigl(\ind[rw<A]+\ind[rw>D]\bigr).
\]

The correspondence is bijective, since $k'=V/w$ is determined by $u,V,w$;
this remains true when $w^2\mid V$. Put
\begin{align*}
F_1(u)&=\min\bigl(p_2-1,\left\lceil A/r\right\rceil-1\bigr),\\
F_2(u)&=p_2-1,\\
F_3(u)&=\min(p_2-1,\lfloor D/r\rfloor).
\end{align*}
In $F_1$, $\lceil\cdot\rceil$ is the ceiling function, and $1$ is subtracted after rounding up.
Set these to zero where $a_u=0$. All three are integers in $[0,X]$, and
\begin{equation}\label{eq:R3-cutoffs}
\vartheta(u,w)=\ind[w\le F_1(u)]+\ind[w\le F_2(u)]
-\ind[w\le F_3(u)].
\end{equation}
Separate the roughness condition by \eqref{eq:separate-roughness} and
each comparison by \eqref{eq:separate-cutoff}.
For fixed $g,g',i$, the two coefficients are
\[
a_u e\Big(\frac{-gP^-(u)+g'F_i(u)}K\Big),\quad
b_V(g,g'):=e\Big(\frac{gP^-(V)}K\Big)
\sum_{\substack{w\mid V,\ w\text{ prime},\ w\ge z}}
e\Big(\frac{-g'w}K\Big).
\]
Their moduli are at most $1$ and $5$, since $V\le X/(4A)<z^6$.
With two Fourier norms and $O(L)$ dyadic intervals,
\eqref{eq:type-II-size} now gives
$O(X^{31/40+\eta/2}L^{7/2})$ for \eqref{eq:R3-stopped}.

\noindent\textit{The families $B_1,B_2,B_3,B_4$.}
Each has the form
\[
\sum_{m\in J}B_j(m)e(\alpha m)
=\sum_{A\le u\le D}\ind_{E_j}(u)
\sum_{v:\ uv\in J}\psi(v,P^-(u))e(\alpha uv).
\]
The ordered prime factorisation of $u$ determines the prefix conditions.
For $B_1$, the condition $p_1^2\le m$ follows from $v\ge p_1$.
The residual satisfies $v>X/(8D)>2$, and one Fourier separation with
$O(L)$ dyadic intervals gives
$O(X^{31/40+\eta/2}L^{5/2})$.

All Type I outer variables are below $X^{27/40}$.
Both variables in every Type II piece are $\ll D\le X/P$, and the
underlying products place them in $[A/8,D]$ up to endpoint constants.
Adding the eight estimates proves the proposition.
\end{proof}

\begin{corollary}\label{cor:minor-lambda}
Uniformly for $\alpha\in\m$,
\[
\sup_{\alpha\in\m}|S_\lambda(\alpha)|
\ll X^{31/40+\eta/2}L^{9/2}=X^{1-t/2}L^{9/2}.
\]
\end{corollary}

\begin{proof}
Partial summation gives
\[
S_\lambda(\alpha)=\log\frac X4\,T_\rho(\alpha;I)
-\int_{X/8}^{X/4}T_\rho\big(\alpha;(X/8,y]\big)\frac{dy}y.
\]
Its modulus is at most $2L$ times the largest prefix sum, and
Proposition~\ref{prop:minor-rho} applies.
\end{proof}

\section{Bilinear distribution with both maxima}\label{sec:bilinear}

\begin{bilinearlemma}\label{lem:B}
Let $X$ be large, $L=\log X$, $z=X^{1/10}$, and $1\le Q<X^{1/2}$.
Let $M,N\ge1$. Suppose that $\alpha,\beta$ satisfy the following conditions.
Their supports lie in $(M,2M]$ and $(N,2N]$, respectively, and consist
of $z$-rough integers greater than $1$.
Their coefficients have modulus at most $c_0$.
For some $\mathcal A\ge1$, every $d\le L^{\mathcal A}$,
$(b,d)=1$ and real $u$ satisfy
\begin{equation}\label{eq:SW-hypothesis}
\left|\sum_{\substack{n\le u\\n\equiv b\pmod d}}\beta_n
-\frac1{\varphi(d)}\sum_{\substack{n\le u\\(n,d)=1}}\beta_n\right|
\le C_{\mathrm{SW}}N L^{-\mathcal A}.
\end{equation}
Put $f=\alpha*\beta$. For $1\le\mathcal D\le\mathcal A$,
\begin{equation}\label{eq:lemma-B}
\begin{split}
\sum_{q\le Q}\Delta_f(q)\ll{}&XL^{2\mathcal D+1-\mathcal A}
+X^{9/10}L^{\mathcal D+1}+XL^{2-\mathcal D}\\
&+X^{19/20}L^3+QX^{1/2}L^2.
\end{split}
\end{equation}
The implied constant depends only on $c_0,C_{\mathrm{SW}}$.
In particular, for $Q\le Q_B$ and $B'>0$, take
$\mathcal D=B'+2$ and $\mathcal A=3B'+5$.
Then the left side is $O(XL^{-B'})$.
\end{bilinearlemma}

\begin{fact}\label{fact:Perron}
\cite[Chapter~17]{Da} For $x>0$, $x\ne1$, $c>0$ and $T\ge1$,
\[
\left|\frac{1}{2\pi i}\int_{c-iT}^{c+iT}\frac{x^s}{s}\,ds
-\ind[x>1]\right|
\ll x^c\min\left(1,\frac{1}{T|\log x|}\right).
\]
\end{fact}

\begin{fact}\label{fact:large-sieve}
\cite[Theorem~7.13]{IK} If $c_n$ is supported on $N'$ consecutive integers,
then
\[
\sum_{r\le R'}\frac r{\varphi(r)}
\sum_{\chi\bmod r}^{*}
\left|\sum_n c_n\chi(n)\right|^2
\le (N'+R'^2)\sum_n|c_n|^2.
\]
The star restricts the sum to primitive characters.
\end{fact}

\begin{proof}[Proof of Lemma~\ref{lem:B}]
We may assume that both sequences are nonzero and $MN<X$, since
otherwise all prefix sums vanish. Roughness gives
\begin{equation}\label{eq:B-support}
z/2\le M,N\le2X/z=2X^{9/10},\qquad
\|\alpha\|_2^2\le2c_0^2M,\quad\|\beta\|_2^2\le2c_0^2N.
\end{equation}
For $\ell\le X$ in the support, $\Omega(\ell)\le10$, hence
$|f(\ell)|\le2^{10}c_0^2$. Write
\[
S_\chi(y;\alpha',\beta')=\sum_{mn\le y}\alpha'_m\beta'_n\chi(mn).
\]
\noindent\textit{Orthogonality and induction.}
The principal character gives the exact coprime mean in $\Delta_f$.
Therefore
\begin{equation}\label{eq:B-orthogonality}
\sum_{q\le Q}\Delta_f(q)
\le\Sigma:=\sum_{q\le Q}\frac1{\varphi(q)}
\sum_{\substack{\chi\bmod q\\\chi\ne\chi_0}}
\max_{y\le X}|S_\chi(y;\alpha,\beta)|.
\end{equation}
Let $\chi$ be induced by a primitive character $\chi_1$ modulo $r$, and
write $q=rk$. For all integers $\ell$,
\[
\chi(\ell)=\chi_1(\ell)\ind_{(\ell,k)=1}.
\]
No coprimality condition between $r,k$ is required. Put $\alpha^{(k)}_m=\alpha_m\ind_{(m,k)=1}$ and
$\beta^{(k)}_n=\beta_n\ind_{(n,k)=1}$.
The induction is a bijection onto the pairs with $r>1$, $rk\le Q$.
Since $\varphi(rk)\ge\varphi(r)\varphi(k)$,
\begin{equation}\label{eq:B-induction}
\Sigma\le\sum_{k\le Q}\frac1{\varphi(k)}
\sum_{1<r\le Q/k}\frac1{\varphi(r)}\sums_{\chi_1\bmod r}
\max_{y\le X}|S_{\chi_1}(y;\alpha^{(k)},\beta^{(k)})|.
\end{equation}
We split this sum at $r=L^{\mathcal D}$.

\noindent\textit{Small conductors.}
Only here do we remove the restriction $(mn,k)=1$.
Since $k\le Q<z^5$, at most four primes at least $z$ divide $k$, and for
each such prime at most $X/z$ products $mn\le X$ occur.
Thus, uniformly in all cutoffs,
\[
\left|S_{\chi_1}(y;\alpha^{(k)},\beta^{(k)})
-S_{\chi_1}(y;\alpha,\beta)\right|
\le4\cdot2^{10}c_0^2X/z.
\]
For $1<r\le L^{\mathcal D}$, expand $\beta$ by reduced residue classes.
The primitive character is nonprincipal, so \eqref{eq:SW-hypothesis} gives
\[
\left|\sum_{n\le u}\beta_n\chi_1(n)\right|
\le\varphi(r)C_{\mathrm{SW}}NL^{-\mathcal A}.
\]
Summing over $m$ gives $|S_{\chi_1}(y;\alpha,\beta)|
\le2c_0C_{\mathrm{SW}}rXL^{-\mathcal A}$, and there are at most
$\varphi(r)$ primitive characters. Using $\sum_{k\le Q}\varphi(k)^{-1}\ll L$, we obtain
\begin{equation}\label{eq:B-small}
\Sigma_{\mathrm{small}}
\ll XL^{2\mathcal D+1-\mathcal A}+X^{9/10}L^{\mathcal D+1}.
\end{equation}

\noindent\textit{Perron's formula and the maximum in $y$.}
Put $c=1/L$ and $T=X^3$. Every prefix can be represented with a
half-integer $y=j+1/2$. The convolution is supported in $(MN,4MN]\subset[1,4X]$,
and its absolute mass is at most $4c_0^2X$.
For the relevant $\ell$, $|\log(y/\ell)|\ge1/(10X)$ and $(y/\ell)^c\le3$,
so (\ref{fact:Perron}) gives a uniform error $O(X^{-1})$. Define
\[
A_k(s,\chi)=\sum_m\alpha_m^{(k)}\chi(m)m^{-s},\qquad
B_k(s,\chi)=\sum_n\beta_n^{(k)}\chi(n)n^{-s}.
\]
Then
\[
\max_{y\le X}|S_\chi(y;\alpha^{(k)},\beta^{(k)})|
\le\frac{3}{2\pi}\int_{-T}^T
\frac{|A_kB_k(c+it,\chi)|}{|c+it|}\,dt+O(X^{-1}).
\]
Since $\int_{-T}^T|c+it|^{-1}dt\le8L$ for large $X$,
any finite character family $\mathcal F$ satisfies
\begin{equation}\label{eq:B-Perron}
\sum_{\chi\in\mathcal F}\max_y|S_\chi(y)|
\le4L\sup_{|t|\le T}\sum_{\chi\in\mathcal F}|A_kB_k(c+it,\chi)|
+O(|\mathcal F|/X).
\end{equation}
The same argument applies with positive character weights, and it keeps
the maximum inside the sum over characters.

\noindent\textit{Large conductors.}
The intervals $(R_i,2R_i]$, where $R_i=2^iL^{\mathcal D}$, $i\ge0$,
cover every conductor greater than $L^{\mathcal D}$.
Fix $k$, one such interval, and $t$.
Cauchy--Schwarz with weights $r/\varphi(r)$, followed by
(\ref{fact:large-sieve}), gives
\begin{align*}
\sum_{R<r\le\min(2R,Q/k)}\frac1{\varphi(r)}
\sum_{\chi_1\bmod r}^{*}|A_kB_k|
&\le\frac1R\big((2M+4R^2)2c_0^2M\big)^{1/2}\\
&\qquad\times\big((2N+4R^2)2c_0^2N\big)^{1/2}\\
&\ll\Big(\frac{X^{1/2}}R+M^{1/2}+N^{1/2}+R\Big)X^{1/2}.
\end{align*}
The weighted Perron error is $O(R/X)$.
Sum \eqref{eq:B-Perron} over $k$ and these intervals; the endpoint terms
form geometric sums. Using \eqref{eq:B-support}, we get
\begin{align}
\Sigma_{\mathrm{large}}
&\ll L^2\Big(XL^{-\mathcal D}
+LX^{1/2}(2X^{9/10})^{1/2}+QX^{1/2}\Big)\notag\\
&\ll XL^{2-\mathcal D}+X^{19/20}L^3+QX^{1/2}L^2.
\label{eq:B-large}
\end{align}
Equations~\eqref{eq:B-small} and \eqref{eq:B-large} prove \eqref{eq:lemma-B}.
With the stated choices of $\mathcal D,\mathcal A$, both logarithmic terms
are $XL^{-B'}$, and the remaining three terms have a fixed power saving.
\end{proof}

\subsection{Siegel--Walfisz estimates for prime boxes}

\begin{fact}\label{fact:SW}
\cite[Chapter~22]{Da} For every $\mathcal A>0$, there is $c_{\mathcal A}>0$
such that
\[
\left|\pi(w;d,b)-\frac{\pi(w)}{\varphi(d)}\right|
\ll_{\mathcal A}w\exp\big(-c_{\mathcal A}\sqrt{\log w}\big),
\]
uniformly for $(b,d)=1$ and $d\le(\log w)^{\mathcal A+1}$.
\end{fact}

\begin{primeboxlemma}\label{lem:C}
Let $\mathcal A>0$.
For $\beta=\ind_{\PP\cap[U,V)}$, where $z\le U<V\le4U$,
the left side of \eqref{eq:SW-hypothesis} is
$O_{\mathcal A}(Ue^{-c_{\mathcal A}\sqrt L})$,
uniformly for $d\le L^{\mathcal A}$, $(b,d)=1$ and real $u$.

Suppose $\beta_1$ satisfies \eqref{eq:SW-hypothesis} with scale $N_1$
and constant $C_1$.
If $\gamma$ is supported on $(N_2,2^jN_2]$ and $|\gamma|\le c$,
then $\beta_1*\gamma$ satisfies it with scale $N_1N_2$
and constant $C_1c2^{j+1}$.
Restricting any such sequence to $(N',2N']$ multiplies its constant by at most $2$.
\end{primeboxlemma}

\begin{proof}
For large $X$, $d\le L^{\mathcal A}<z$, so every prime in $[U,V)$ is coprime to $d$.
Put $E(w)=\pi(w;d,b)-\pi(w)/\varphi(d)$.
The discrepancy is zero for $u<U$. For $U\le u<V$ it equals
$E(u)-E(U^-)$, and for $u\ge V$ it equals $E(V^-)-E(U^-)$; the left limits preserve
both half-open endpoints. All arguments satisfy $\log w\ge L/10$, so for large $X$ we have
$d\le(\log w)^{\mathcal A+1}$, and (\ref{fact:SW}) applies.

For the convolution, fix $n_2$ in the support of $\gamma$.
If $(n_2,d)>1$, neither the reduced residue sum nor its coprime mean contributes.
Otherwise the residue for $n_1$ is $bn_2^{-1}\pmod d$, and we apply
\eqref{eq:SW-hypothesis} at the cutoff $u/n_2$. The errors sum to at most
\[
C_1N_1L^{-\mathcal A}\sum_{n_2}|\gamma(n_2)|
\le C_1c\,2^{j+1}N_1N_2L^{-\mathcal A}.
\]
The main terms give the coprime mean of the complete convolution.
Finally, a dyadic restriction is the difference of two prefix sums.
\end{proof}

\section{Distribution of the signed weight}\label{sec:BV9}

\begin{proposition}[Distribution of the signed weight]\label{prop:BV9}
For every $B>0$,
\begin{equation}\label{eq:BV9}
\sum_{d\le Q_B}2^{\omega(d)}\Delta_\lambda(d)
\ll_{B,\eta}XL^{-B},\qquad Q_B=X^{9/20-\eta}.
\end{equation}
Both maxima in \eqref{eq:Delta} remain inside the modulus sum.
This estimate is unconditional.
\end{proposition}

\subsection{Marked prefixes on the full support}

Extend each discarded family to all positive integers $m$. Use the
prefix conditions in Table~\ref{tab:families}, including $p_1^2\le m$,
and the residual $\psi(m/P_j,p_j)$. Write their sum as $\widetilde U$,
and put
\[
\widetilde U^\flat(m):=\widetilde U(m)\ind_{m\le X/4}.
\]
On $I$ this equals the sum of the four losses. It is supported on $z$-rough integers with $\Omega(m)\le9$, and
\begin{equation}\label{eq:Uflat-bound}
0\le\widetilde U^\flat(m)\le\binom92+\binom94=162.
\end{equation}
For each four-prime prefix, the three possibilities $L_4,U_4,V_4$ are disjoint.

Let $m=\ell_1\cdots\ell_r$ be squarefree, with
$\ell_1>\cdots>\ell_r\ge z$.
Put $t_i=\log\ell_i/L$ and $s=\sum_i t_i$.
A marked prefix has residual roughness at its smallest selected prime.
This holds precisely when that prime is $\ell_r$, also when the residual
is $1$. Thus $\widetilde U(m)=G_r(t)$, where
\begin{equation}\label{eq:marked-prefixes}
\begin{split}
G_r(t)={}&\sum_{i<r}
\ind[t_i<9/20,\ t_i+t_r>11/20,\ 2t_i\le s]\\
&+\sum_{i_1<i_2<i_3<r}
\ind[t_{i_1}+t_{i_2}+t_{i_3}+t_r<9/20,\ 2t_{i_1}\le s]\\
&+\sum_{i_1<i_2<i_3<r}
\ind[t_{i_1}+t_{i_2}+t_{i_3}<9/20,\\
&\hspace{24mm}t_{i_1}+t_{i_2}+t_{i_3}+t_r>11/20,\ 2t_{i_1}\le s]\\
&+\sum_{i_1<i_2<i_3<r}
\ind[t_{i_1}+t_{i_2}<9/20,
t_{i_1}+t_{i_2}+t_{i_3}>11/20,\ 2t_{i_1}\le s].
\end{split}
\end{equation}
In particular, $G_1=0$ and $\widetilde U(1)=0$.
For $2\le r\le9$, let $\mathcal H_r$ consist of the hyperplanes
\begin{equation}\label{eq:boundary-hyperplanes}
\sum_{i\in J}t_i=9/20\ \text{or}\ 11/20\quad(1\le|J|\le4),
\quad 2t_i=\sum_kt_k,\quad t_i=t_j\quad(i\ne j).
\end{equation}
These form a fixed finite rational arrangement. Off it, $G_r$ is constant
in each component of the strictly ordered cone, and the bound $162$ holds
there for $r\le9$.

\subsection{Initial reductions and single primes}

Partial summation and \eqref{eq:cutoffs} give
\begin{equation}\label{eq:BV-reductions}
\Delta_\lambda(d)\le2L\Delta_\rho(d),\qquad
\Delta_\rho(d)\le
\Delta_{\ind_\PP\ind_I}(d)+\Delta_{\widetilde U^\flat\ind_I}(d).
\end{equation}
We split $\widetilde U^\flat$ into squarefree and nonsquarefree parts
and add the divisor weight last.

\begin{fact}\label{fact:classical-BV}
\cite[Chapter~28]{Da} For every $H>0$, some $K(H)>0$ satisfies
\begin{equation}
\sum_{d\le X^{1/2}L^{-K(H)}}
\max_{(a,d)=1}\max_{y\le X}
\left|\sum_{\substack{n\le y\\n\equiv a\pmod d}}\Lambda(n)
-\frac{y}{\varphi(d)}\right|
\ll_H XL^{-H}.
\end{equation}
\end{fact}

Since $Q_B/(X^{1/2}L^{-K(H)})=X^{-1/20-\eta}L^{K(H)}\to0$,
this covers all required moduli. Prime powers contribute
$O(\sqrt X L)$ per modulus. Partial summation converts prime logarithmic
sums to prime counts and retains the maximum in $y$. The modulus
$1$ case controls the corresponding total prime count.
Removing primes that divide $d$ costs $O(\omega(d)/\varphi(d))$, so the
three conversions have total cost
\[
Q_BX^{1/2}L+\sum_{d\le Q_B}\frac{\omega(d)}{\varphi(d)}
\ll X^{19/20-\eta}L.
\]
Consequently, for every $H>0$,
\begin{equation}\label{eq:prime-BV}
\sum_{d\le Q_B}\Delta_{\ind_\PP}(d)\ll_H XL^{-H}.
\end{equation}
Single primes are handled by this estimate; Lemma~\ref{lem:B} is not
used for them.

\subsection{Absolute-coordinate cubes}

Fix $C_0>1$. Put
\[
\varpi:=\frac1{10N_\varpi},\qquad
N_\varpi=\lceil L^{C_0}\rceil,\qquad
\frac1{20}L^{-C_0}\le\varpi\le\frac1{10}L^{-C_0},
\qquad 1/10\in\varpi\mathbb Z.
\]
For $2\le r\le9$, use cubes
\begin{equation}\label{eq:cubes}
C=\prod_{i=1}^r[a_i,a_i+\varpi),\qquad
a_i\in\tfrac1{10}+\varpi\Z_{\ge0},\quad
a_i<1,\quad a_1\ge\cdots\ge a_r.
\end{equation}
Let $\mathcal B_i(C)=\PP\cap[X^{a_i},X^{a_i+\varpi})$ and
\[
\conv_C:=\ind_{\mathcal B_1(C)}*\cdots*\ind_{\mathcal B_r(C)}.
\]
Every prime is at least $z$, the box ratio is
$X^\varpi\le\exp(L^{1-C_0}/10)\le2$, and there are $O(L^{9C_0})$ cubes.
An interior cube has closure disjoint from all hyperplanes in
$\mathcal H_r$, and every other cube is a boundary cube.
Only dimensions $2\le r\le9$ are used, since the discarded weight vanishes
when $r\le1$.

For an interior cube, the prime boxes are pairwise disjoint and ordered,
so unique factorisation gives $\conv_C(m)\in\{0,1\}$. It is $1$ precisely for the squarefree integers with $\Omega(m)=r$ and $t(m)\in C$.
Let $c_C\in[0,162]\cap\Z$ be the constant value of $G_r$ on $C$. Then,
for every integer $m\ge1$,
\begin{equation}\label{eq:interior-identity}
\widetilde U^\flat(m)\ind_{\mathrm{sqf}}(m)\ind_{t(m)\in C}
=c_C\conv_C(m)\ind_{m\le X/4}.
\end{equation}
The cutoff on the right is retained, and $c_C$ counts all admissible
marked prefixes.

\subsection{Distribution of each convolution}

For any cube, including boundary cubes, split
\[
\mathrm{conv}_C=\alpha*\beta,\qquad
\alpha=\ind_{\mathcal B_1(C)},\qquad
\beta=\ind_{\mathcal B_2(C)}*\cdots*\ind_{\mathcal B_r(C)}.
\]
Both sequences are supported on rough integers greater than $1$, and
their coefficients are bounded by $1$ and $(r-1)!\le8!$.
Lemma~\ref{lem:C} supplies \eqref{eq:SW-hypothesis} at every fixed level.
For $r=2$, use its prime-box assertion. For $r>2$, convolve one prime box
with the other $r-2$ boxes; the latter convolution has coefficients at
most $7!$ and bounded support ratio.
All constants are uniform over the cubes.

We specify the dyadic endpoints.
Put $N_\star=\prod_{i=2}^rX^{a_i}$. The support of $\beta$ lies in
$[N_\star,2^{r-1}N_\star)$, and we use the $r$ intervals
\[
(2^{j-1}N_\star,2^jN_\star],\qquad 0\le j\le r-1.
\]
They begin at $N_\star/2$, so no integer at $N_\star$ is omitted.
Use two intervals anchored at $X^{a_1}/2$ for $\alpha$.
Lemma~\ref{lem:C} controls each restriction, and changing to the dyadic
scale costs at most a factor $2^9$.
Lemma~\ref{lem:B} therefore gives, uniformly in $C$,
\begin{equation}\label{eq:cube-BV}
\sum_{d\le Q_B}\Delta_{\conv_C\ind_{[1,X/4]}}(d)
\le\sum_{d\le Q_B}\Delta_{\conv_C}(d)\ll_{B'}XL^{-B'}.
\end{equation}
The distribution estimate is applied to the full product $m$ before
truncation, and the modulus remains $d$ throughout.

\subsection{Boundary cubes and positive transfer}

Define
\begin{gather*}
\mathcal E(m):=\sum_{C\text{ boundary}}
\widetilde U^\flat(m)\ind_{\mathrm{sqf}}(m)\ind_{t(m)\in C},\\
W:=\sum_{C\text{ boundary}}\mathrm{conv}_C\ge0,
\qquad W^\flat:=W\ind_{[1,X/4]}.
\end{gather*}
Then $W\ge0$ and $|\mathcal E|\le162W^\flat$. Every relevant sorted
vector lies in exactly one cube. Hence
\begin{equation}\label{eq:squarefree-decomposition}
\widetilde U^\flat\ind_{\mathrm{sqf}}
=\sum_{C\ \mathrm{interior}}c_C\conv_C\ind_{[1,X/4]}+\mathcal E.
\end{equation}

\begin{lemma}\label{lem:positive-transfer}
If $|E|\le\mathcal CW$ and $W\ge0$, then
\[
\Delta_{E}(d)
\le\mathcal C\Delta_W(d)
+\frac{2\mathcal C}{\varphi(d)}\sum_{m\le X}W(m).
\]
\end{lemma}

\begin{proof}
Bound each of the two terms defining the discrepancy of $E$ by its
absolute value, and replace $|E|$ by $\mathcal CW$.
The residue sum of $W$ is at most its coprime mean plus $\Delta_W(d)$,
and the two means are each bounded by
$\varphi(d)^{-1}\sum_{m\le X}W(m)$. These bounds are uniform in the
residue and the cutoff.
\end{proof}

By \eqref{eq:cube-BV},
\begin{equation}\label{eq:boundary-BV}
\sum_{d\le Q_B}\Delta_{W^\flat}(d)
\ll L^{9C_0}XL^{-B'}.
\end{equation}
This includes cubes meeting the diagonals. For a term with $r\le9$ prime
boxes, split off one box as $\alpha$ and put the remaining $r-1$ boxes in
$\beta$. The $\beta$ coefficients are at most $(r-1)!\le8!$, which is the
bound used by Lemma~\ref{lem:B}. The full convolution can have coefficients
up to $r!\le9!$. There are no one-dimensional boundary convolutions.

\begin{lemma}\label{lem:boundary-mass}
We have
\[
\sum_{m\le X}W(m)
\ll\varpi LX+X/z
\ll XL^{1-C_0}+X^{9/10}.
\]
\end{lemma}

\begin{proof}
A tuple in a boundary cube lies within distance $\varpi$, in the maximum norm,
of some hyperplane $c\cdot t=c_0$ in \eqref{eq:boundary-hyperplanes}.
Choose $i_0$ with $c_{i_0}\ne0$, and fix the other primes, whose product
is $M'\le X/z$. The remaining prime lies in
\[
[T,Te^{\kappa\varpi L}]\cap[z,X/M']
\]
This interval contains $O(\varpi LX/M'+1)$ integers, and by
(\ref{fact:PNT}),
\[
\sum_{\text{tuples}}\frac1{M'}\le\left(\sum_{z\le p\le X}\frac1p\right)^{r-1}\le3^8
\]
for large $X$.
There are at most $(r-1)!\lfloor X/z\rfloor\le8!X/z$ labelled tuples
with product at most $X/z$. The number of hyperplanes and dimensions is
fixed, so summing proves the assertion.
\end{proof}

Using $\sum_{d\le Q_B}\varphi(d)^{-1}\ll L$, the last two lemmas give
\begin{equation}\label{eq:squarefree-BV}
\sum_{d\le Q_B}\Delta_{\widetilde U^\flat\ind_{\mathrm{sqf}}}(d)
\ll L^{9C_0}XL^{-B'}+L(XL^{1-C_0}+X^{9/10}).
\end{equation}

\subsection{Repeated prime factors}

The nonsquarefree part has modulus at most $162$, and it is supported on
$p^2\mid m$ for some prime $p\ge z$. Fix $(a,d)=1$.
If $p\mid d$, this divisibility is impossible in the reduced residue
class. If $p\nmid d$, the Chinese remainder theorem gives
\[
\#\{m\le X:m\equiv a\pmod d,\ p^2\mid m\}
\le X/(dp^2)+1.
\]
Summing over $z\le p\le\sqrt X$ gives $O(X/(dz)+\sqrt X)$, and the
coprime-average term is $O(X/(\varphi(d)z))$. Therefore
\begin{equation}\label{eq:nonsquarefree-BV}
\sum_{d\le Q_B}\Delta_{\widetilde U^\flat\ind_{\mathrm{nsqf}}}(d)
\ll XL/z+Q_B\sqrt X
=X^{9/10}L+X^{19/20-\eta}.
\end{equation}
Repeated factors remain in the weight. This step counts multiples
directly and needs no distribution estimate at modulus $dp^2$.

\subsection{The divisor weight and the parameter choices}

Since $|\lambda|\le256L$,
\[
\Delta_\lambda(d)\ll XL/\varphi(d)+L.
\]
The elementary divisor estimates
\[
\sum_{d\le Q}\frac{4^{\omega(d)}}{\varphi(d)}\ll L^4,
\qquad
\sum_{d\le Q}4^{\omega(d)}\ll QL^3
\]
give
\begin{equation}\label{eq:four-omega}
\sum_{d\le Q_B}4^{\omega(d)}\Delta_\lambda(d)\ll XL^5.
\end{equation}
To prove them, expand $4^{\omega(d)}$ over squarefree divisors and use
$\varphi(ab)\ge\varphi(a)\varphi(b)$; the corresponding Euler products
have logarithmic exponents $3$ and $4$. Cauchy--Schwarz yields
\begin{equation}\label{eq:divisor-CS}
\sum_{d\le Q_B}2^{\omega(d)}\Delta_\lambda(d)
\le\left(\sum_{d\le Q_B}\Delta_\lambda(d)\right)^{1/2}
\left(\sum_{d\le Q_B}4^{\omega(d)}\Delta_\lambda(d)\right)^{1/2}.
\end{equation}

\begin{proof}[Proof of Proposition~\ref{prop:BV9}]
Given $B>0$, choose, in this order,
\[
H:=2B+8,\qquad C_0:=H+3,\qquad B':=H+9C_0+1,
\]
then $\mathcal D=B'+2$ and $\mathcal A=3B'+5$, and finally take $X$
sufficiently large.
Equations~\eqref{eq:prime-BV}, \eqref{eq:squarefree-BV},
\eqref{eq:nonsquarefree-BV} and \eqref{eq:BV-reductions} give
\[
\sum_{d\le Q_B}\Delta_\rho(d)\ll XL^{-H},
\qquad
\sum_{d\le Q_B}\Delta_\lambda(d)\ll XL^{1-H}.
\]
Combining \eqref{eq:four-omega} and \eqref{eq:divisor-CS} gives
$XL^{3-H/2}=XL^{-B-1}$, which proves \eqref{eq:BV9}.
\end{proof}

\section{Major arcs}\label{sec:major}

Only this section uses $H_{7/8}$. We first prove the prime-side mean
square. The argument follows \cite[Section~8]{Pi}, with the density exponents computed below.

\subsection{Classical inputs}

Let $N^*(\sigma,T,Q')$ count, with multiplicity, the zeros
$\varrho=\beta+i\gamma$ satisfying $\beta\ge\sigma$ and $|\gamma|\le T$, summed over primitive
characters of conductor at most $Q'$.

\begin{fact}\label{fact:HB-one}
\cite[Theorem~1, p.~231]{HB} For $11/14\le\sigma\le1$,
\[
N^*(\sigma,T,Q')\ll_\varepsilon(Q'^2T)^{2(1-\sigma)+\varepsilon}.
\]

\end{fact}

\begin{fact}\label{fact:HB-twenty}
\cite[Theorem~2(3), p.~232]{HB}; \cite[Lemma~4.17]{Pi}.
Uniformly for $1/2\le\sigma\le1$,
\[
N^*(\sigma,T,Q')\ll_\varepsilon
(Q'^2T^{6/5})^{(20/9)(1-\sigma)+\varepsilon}.
\]

\end{fact}

\begin{fact}\label{fact:Montgomery}
\cite[Theorem~12.2]{C2c}; \cite[Lemma~4.16]{Pi}.
For $1/2\le\sigma\le1$,
\[
N^*(\sigma,T,Q')\ll
(Q'^2T)^{3(1-\sigma)/(2-\sigma)}\log^9(Q'T).
\]
\end{fact}

\begin{fact}\label{fact:Gallagher}
\cite[Lemma~1]{C4}; \cite[Lemma~4.11]{Pi}.
For a finite sequence $a_n$ and $\kappa>0$,
\[
\int_{-\kappa}^{\kappa}
\left|\sum_n a_ne(n\beta)\right|^2d\beta
\ll \kappa^2\int
\left|\sum_{x<n\le x+1/(2\kappa)}a_n\right|^2dx.
\]

\end{fact}

\begin{fact}\label{fact:Gauss}
\cite[Section~3.4]{IK}; \cite[Chapter~9]{Da}.
If $\chi\bmod q$ is induced by a primitive character of conductor $r$, then
$|\tau(\chi)|^2\le r$. Also $\tau(\chi_{0,q})=\mu(q)$.
\end{fact}

\begin{fact}\label{fact:explicit}
\cite[Chapters~15--19]{Da}; \cite[Lemma~4.15]{Pi}.
The explicit formula for primitive Dirichlet characters has the uniform
truncation error
\[
O\!\left(\frac{x\log^2(rxT)}{T}+x^{1/4}\log x\right).
\]
The number of zeros in a unit height interval is
$O(\log(r(|T|+2)))$.
\end{fact}

\subsection{The prime-side mean square}

On $\M(q,a)$, put
\begin{equation}\label{eq:S1-definition}
\mathcal T(\beta)=\sum_{X_1<k\le X}e(k\beta),\qquad
S_1(a/q+\beta)=S(a/q+\beta)-\frac{\mu(q)}{\varphi(q)}\mathcal T(\beta).
\end{equation}
The arcs are disjoint, so this defines $S_1$ on $\M$.

\begin{proposition}\label{prop:prime-mean-square}
Assume $H_{7/8}$. Let $0<\eta\le1/100$, $0<\varepsilon_0\le\eta/10$
and $X_1=X^{1-\varepsilon_0}$.
Uniformly for $2\le P\le X^{9/20-\eta}$,
\begin{equation}\label{eq:prime-mean-square}
\int_\M|S_1(\alpha)|^2d\alpha
\ll_\eta L^{18}X^{1-8\eta/9}+L^8P^2
\ll_\eta L^{18}X^{1-8\eta/9}.
\end{equation}
\end{proposition}

\begin{proof}
Put $\vartheta=\log P/L$, $Q=X/P$ and $\epsilon=\eta/90$, and write
$\varrho=1-\delta+i\gamma$.

\noindent\textit{Character expansion.}
Every prime in $S$ is coprime to $q\le P$, since $P<\sqrt X<X_1$.
For the character $\chi^*$ inducing $\chi\bmod q$,
\[
S(a/q+\beta)=\frac1{\varphi(q)}\sum_{\chi\bmod q}
\chi(a)\tau(\bar\chi)
\sum_{X_1<p\le X}\chi^*(p)\log p\,e(\beta p).
\]
There is one exterior factor $1/\varphi(q)$. For primitive $\chi\bmod r$, put
\[
\begin{aligned}
S_1(\chi^*,\beta)&=\sum_n a_n(\chi^*)e(n\beta),\\
a_n(\chi^*)&=\ind_{X_1<n\le X}
\bigl(\ind_{\PP}(n)\chi^*(n)\log n-E(\chi^*)\bigr),\\
E(\chi^*)&=\ind[r=1].
\end{aligned}
\]

Orthogonality in $a$, followed by (\ref{fact:Gauss}), gives
\begin{equation}\label{eq:character-mean-square}
\begin{gathered}
\int_\M|S_1|^2\ll L^2\sum_{r\le P}\;\sums_{\chi\bmod r}W_1(\chi)^2,
\\
W_1(\chi)^2=\int_{|\beta|\le1/(rQ)}
\left|\sum_na_n(\chi)e(n\beta)\right|^2d\beta.
\end{gathered}
\end{equation}
Here we used
\[
\sum_{\substack{q\le P\\r\mid q}}\frac1{\varphi(q)}
\le \frac1{\varphi(r)}\sum_{l\le P/r}\frac1{\varphi(l)}
\ll \frac L{\varphi(r)},
\qquad \frac r{\varphi(r)}\ll L.
\]
This includes all moduli $q\le P$.

\noindent\textit{Short intervals.}
Apply (\ref{fact:Gallagher}) with $Y=rQ/2\le X/2$, and let
$X_2=\max(X_1,6Y)$. Then
\begin{equation}\label{eq:Gallagher-ranges}
W_1(\chi)^2\ll\int_{X_1-Y}^X
\left|Y^{-1}\sum_{x<n\le x+Y}a_n(\chi)\right|^2dx
=I_1+I_2+I_3.
\end{equation}
With empty intervals omitted, the three ranges are
\[
I_1=[X_1-Y,\min(X_2,X-Y)],\qquad
I_2=[X_2,X-Y],\qquad
I_3=[X-Y,X].
\]
The first has length at most $7Y$.

\noindent\textit{The explicit formula.}
Put $T_0=3\sqrt X$ and $\Lambda_\PP(n)=\ind_\PP(n)\log n$. By (\ref{fact:explicit}), uniformly for $X_1/2\le u_1<u_2\le9X$,
\begin{equation}\label{eq:explicit-difference}
\sum_{u_1<n\le u_2}\bigl(\chi(n)\Lambda_\PP(n)-E(\chi)\bigr)
=-\sum_{\substack{\varrho:\ \beta\ge1/2\\|\gamma|\le T_0}}
\frac{u_2^\varrho-u_1^\varrho}{\varrho}
+O(\sqrt X L^2).
\end{equation}
Here prime powers cost $O(\sqrt XL)$, the zeros with $\beta<1/2$ cost
$O(\sqrt XL^2)$, and the truncation costs $O(XL^2/T_0+X^{1/4}L)$.
The pole cancels the term $E(\chi)$ up to $O(1)$. Constant terms cancel
in the difference, and their logarithmic terms cost $O(L)$.
Hypothesis $H_{7/8}$ and the functional equation keep all nontrivial zeros
in $1/8\le\beta\le7/8$, so every zero retained in \eqref{eq:explicit-difference} has
\begin{equation}\label{eq:delta-range}
1/8\le\delta\le1/2.
\end{equation}
The squared error in the two end ranges contributes at most
\[
\sum_{r\le P}rY^{-1}XL^4\ll L^4P^2.
\]

The middle-range error is treated below.

\noindent\textit{Classes of zeros.}
Partition conductors into intervals $R<r\le2R$ with dyadic $R$,
including the initial interval $R=1/2$ to cover $r=1$. For $\mathcal M=2^\mu$, $\mu\ge0$, use the height intervals
\[
(\mathcal M-1)P/R\le |\gamma|\le(2\mathcal M-1)P/R.
\]

Shared endpoints are assigned once.
Also use $h-1/L<\delta\le h$, with $h\in L^{-1}\Z$.
There are $O(L^3)$ classes, and $O(L^2)$ for each character. By \eqref{eq:delta-range}, only $1/8\le h\le1/2+1/L$ occur.
Within a class,
\[
X^{-\delta}\le eX^{-h},\qquad {X_1}^{-\delta}\le e\,{X_1}^{-h},
\qquad |\gamma|\le T_{\mathcal M}:=2\mathcal MP/R,
\qquad T_{\mathcal M}\ge2.
\]
For one character the class contains $O((\mathcal MP/R)L)$ zeros, and
over all characters it contains at most
$N^*(\sigma_h,T_{\mathcal M},2R)$, where
\begin{equation}\label{eq:class-K}
\sigma_h=\max(1-h,1/2),\qquad
\mathcal K(R,\mathcal M,h)
=\mathcal M^{-1}N^*(\sigma_h,2\mathcal MP/R,2R)\,{X_1}^{-2h}.
\end{equation}
For the terminal classes $h>1/2$, the value $\sigma_h=1/2$ is valid
because \eqref{eq:explicit-difference} already restricts to $\beta\ge1/2$.
We also have
\begin{equation}\label{eq:height-decay}
\min\left(1,\frac{2P}{r|\gamma|}\right)\le\frac4{\mathcal M}.
\end{equation}

\noindent\textit{The upper end range.}
For $x\in[X-Y,X]$, the zero term is bounded by
\[
\left|\frac{X^{\varrho}-x^{\varrho}}{Y\varrho}\right|
\le\min\left(x^{-\delta},\frac{2X^{1-\delta}}{Y|\gamma|}\right)
\ll X^{-\delta}\min\left(1,\frac{2P}{r|\gamma|}\right)
\ll\frac{X^{-h}}{\mathcal M}.
\]
Apply Cauchy--Schwarz first to the classes for one character and then
to the zeros in one class. For fixed $R,\mathcal M,h$, its contribution after summing characters is
\begin{multline*}
L^2\cdot RQ\cdot\frac{X^{-2h}}{\mathcal M^2}
\cdot\frac{\mathcal MP}{R}L\cdot
N^*(\sigma_h,T_{\mathcal M},2R)\\
=L^3X\mathcal M^{-1}N^*(\sigma_h,T_{\mathcal M},2R)X^{-2h}
\le L^4X\mathcal K(R,\mathcal M,h).
\end{multline*}
Summing all classes gives $O(L^7X\max\mathcal K)$.

\noindent\textit{The lower end range.}
Here the window is $(\max(x,X_1),x+Y]$, and its zero term is at most
\begin{multline*}
\left|
\frac{(x+Y)^{\varrho}-\max(x,X_1)^{\varrho}}{Y\varrho}
\right|
\le\min\left({X_1}^{-\delta},\frac{2X^{1-\delta}}{Y|\gamma|}\right)\\
\le {X_1}^{-\delta}\min\left(1,\frac{2X}{Y|\gamma|}\right)
\ll\frac{{X_1}^{-h}}{\mathcal M}.
\end{multline*}
Here ${X_1}^{-\delta}$ is a power of $X_1=X^{1-\varepsilon_0}$ with exponent $-\delta$.
The numerator $2X^{1-\delta}$ is a power of $X$ itself, with exponent $1-\delta$.

The same calculation gives $O(L^7X\max\mathcal K)$. This range requires
the factor ${X_1}^{-2h}$ in place of $X^{-2h}$.

\noindent\textit{The middle range.}
Partition it at $x_j=2^jX_2$, using
$x_j,\min(2x_j,X-Y)$. The final interval is truncated at $X-Y$, and
every window remains inside $X_1,X$. Write
\[
f(u)=\sum_{n\le u}\bigl(\chi(n)\Lambda_\PP(n)-E(\chi)\bigr).
\]
For $x\ge X_2\ge6Y$, the following estimate also bounds a truncated interval:
\begin{equation}\label{eq:relative-shift}
\int_x^{2x}\left|\frac{f(u+Y)-f(u)}Y\right|^2du
\le\frac{30}{Y^2}
\max_{Y/(3x)\le\theta'\le3Y/x}
\int_x^{3x}|f(u+\theta'u)-f(u)|^2du.
\end{equation}
To prove it, for $2Y\le v\le3Y$ write
\[
g(u)=f(u+Y)-f(u)
=\bigl[f(u+v)-f(u)\bigr]
-\bigl[f((u+Y)+(v-Y))-f(u+Y)\bigr].
\]
Average over $v$ and use Cauchy--Schwarz.
In the first bracket substitute $\theta'=v/u$; its contribution is at most
$(2/Y)(2x)(3Y/x)=12$ times the maximum integral.
In the second substitute $\theta'=(v-Y)/(u+Y)$; its contribution is at
most $(2/Y)(3x)(3Y/x)=18$ times that maximum.
The common range is $Y/(3x),3Y/x$, and division by $Y^2$ proves
\eqref{eq:relative-shift}.

Here $\theta'\le1$, and by \eqref{eq:explicit-difference} the zero
coefficient is
\begin{gather*}
f(u+\theta'u)-f(u)
=-\sum_{\varrho}c_{\varrho}u^{\varrho}+O(\sqrt X L^2),\\
c_{\varrho}=\frac{(1+\theta')^{\varrho}-1}{\varrho},
\qquad |c_{\varrho}|\le\min\left(\theta',\frac3{|\varrho|}\right).
\end{gather*}
Also
\[
\left|\int_x^{3x}u^{\varrho+\bar\varrho'}\,du\right|
\ll\frac{x^{3-\delta-\delta'}}{1+|\gamma-\gamma'|}.
\]

The unit-interval zero count in (\ref{fact:explicit}) gives
$\sum_{\varrho'}(1+|\gamma-\gamma'|)^{-1}\ll L^2$ for one character.
The arithmetic-geometric mean inequality in each pair yields
\begin{align*}
\frac1{Y^2}\int_x^{3x}\left|\sum_\varrho c_\varrho u^\varrho\right|^2du
&\ll\frac{L^2}{Y^2}\sum_\varrho
x^{3-2\delta}\min(\theta',3/|\varrho|)^2\\
&\ll\frac{L^2}{Y}\sum_\varrho
x^{2-2\delta}\min(\theta',3/|\varrho|).
\end{align*}
Each zero contributes at most
\[
L^2x^{1-2\delta}\min\left(1,\frac{x}{Y|\gamma|}\right)
\ll L^2\frac{X^{1-2h}}{\mathcal M}.
\]
Here $X^{1-2h}=X\cdot X^{-2h}$ is a power of $X$ itself.
Since $X_1<X$, the factor $X^{-2h}$ is at most the separate class factor ${X_1}^{-2h}$ in \eqref{eq:class-K}.
We used $1-2\delta\ge0$ and $Y\ge RX/(2P)$.
Summing the class counts and the $O(L)$ intervals gives the total zero
contribution $O(L^6X\max\mathcal K)$. The error in \eqref{eq:explicit-difference} contributes, per interval,
$O(X^2L^4/Y^2)$, and after the interval and character sums it is
\[
\sum_{r\le P}r\cdot\frac{X^2L^5}{Y^2}
\ll\sum_{r\le P}r\cdot
\frac{X^2L^5}{(rX/(2P))^2}
\ll L^6P^2.
\]
Combining the three ranges gives
\begin{equation}\label{eq:W-class-bound}
\sum_{r\le P}\;\sums_{\chi\bmod r}W_1(\chi)^2
\ll L^7X\max\mathcal K+L^6P^2.
\end{equation}

\noindent\textit{Density estimates for the classes.}
Recall $\vartheta\le9/20-\eta$, $\epsilon=\eta/90$ and
$\varepsilon_0\le\eta/10$.
First suppose $1/8\le h\le3/14$, and use (\ref{fact:HB-one}) with
$Q'=2R$ and $T=2\mathcal MP/R$. Then $Q'^2T\le8P^2\mathcal M$, and the
exponent $2h+\epsilon-1$ of $\mathcal M$ is negative. Thus $\mathcal K\ll X^{-s_3(h)}$, where
\begin{equation}\label{eq:density-s3}
\begin{split}
s_3(h)&=2h(1-\varepsilon_0-2\vartheta)-2\vartheta\epsilon\\
&\ge\frac14\Bigl(\frac1{10}+2\eta-\varepsilon_0\Bigr)
-\frac9{10}\epsilon
\ge\tfrac1{40}+\tfrac{93}{200}\eta\ge\tfrac1{40}.
\end{split}
\end{equation}
The exact positivity condition is
$1-\varepsilon_0-2\vartheta>\vartheta\epsilon/h$, and at $h=1/8$ its
right side is $8\vartheta\epsilon$.

Next suppose $3/14\le h\le3/8-\epsilon$, and use (\ref{fact:HB-twenty}).
Here
\[
Q'^2T^{6/5}=4\cdot2^{6/5}R^{4/5}(\mathcal MP)^{6/5}
\le10P^2\mathcal M^{6/5}.
\]
The exponent of $\mathcal M$ is
$(8/3)h+(6/5)\epsilon-1\le-22\epsilon/15<0$.
Consequently $\mathcal K\ll_\epsilon X^{-s_2(h)}$, with
\begin{equation}\label{eq:density-s2}
\begin{split}
s_2(h)&=2h\Bigl(1-\varepsilon_0-\frac{20}9\vartheta\Bigr)
-2\vartheta\epsilon\\
&\ge\Bigl(\frac37\frac{191}{90}-\frac1{100}\Bigr)\eta\\
&\ge\frac89\eta.
\end{split}
\end{equation}
The required margin is
\begin{equation}\label{eq:critical-density-margin}
1-\varepsilon_0-(20/9)\vartheta\ge191\eta/90>0.
\end{equation}

For $3/8-\epsilon\le h\le1/2$, use (\ref{fact:Montgomery}). The
exponent $3h/(1+h)-1$ of $\mathcal M$ is nonpositive, and we obtain $\mathcal K\ll L^9X^{-s_1(h)}$, where
\begin{equation}\label{eq:density-s1}
\begin{split}
s_1(h)&=\frac{2h}{1+h}
\Bigl(1+h-3\vartheta-\varepsilon_0(1+h)\Bigr)\\
&\ge\frac{2h}{1+h}
\Bigl(\frac1{40}+3\eta-\epsilon-\frac32\varepsilon_0\Bigr)
\ge\frac1{75}.
\end{split}
\end{equation}
The threshold $\vartheta<11/24$ has margin $1/120+\eta$.

Finally, if $1/2<h\le1/2+1/L$, then $\sigma_h=1/2$, and the elementary
zero count gives
\[
N^*(1/2,T,Q')\le N^*(0,T,Q')\ll Q'^2TL=8R\mathcal MP L.
\]
Since $2h>1$,
\[
\mathcal K\ll\mathcal M^{-1}\cdot8R\mathcal MP L^9\cdot {X_1}^{-1}
\le8L^9P^2X^{-(1-\varepsilon_0)}\ll L^9X^{-1/10}.
\]
All classes have now been covered. For $\eta\le1/100$,
\begin{equation}\label{eq:maximum-K}
\max\mathcal K\ll_\eta L^9X^{-c(\eta)},\qquad
c(\eta)=\min\bigl(\tfrac1{40},\tfrac89\eta,\tfrac1{75},\tfrac1{10}\bigr)
=\frac89\eta.
\end{equation}
Equations~\eqref{eq:character-mean-square}, \eqref{eq:W-class-bound}
and \eqref{eq:maximum-K} prove \eqref{eq:prime-mean-square}.
The term $L^8P^2$ is absorbed because $P^2\le X^{9/10-2\eta}$.
\end{proof}

\subsection{A variant without the first density estimate}

One may omit (\ref{fact:HB-one}) and use (\ref{fact:HB-twenty})
throughout $1/8\le h\le3/8-\epsilon$.
The minimum in \eqref{eq:density-s2} is then at $h=1/8$, giving
\[
s_2\ge\frac14\cdot\frac{191}{90}\eta-\frac{\eta}{100}
=0.5206\ldots\eta\ge\frac\eta2.
\]
The mean-square saving is $\eta/2$ in place of $8\eta/9$, and the mixed
error below becomes $X^{1-\eta/4}L^{19/2}$. All other steps are unchanged.

\subsection{The discrete pole and its kernel}

Parseval and Lemma~\ref{lem:weight-norms}, together with
Proposition~\ref{prop:prime-mean-square}, give
\begin{equation}\label{eq:S1-removal}
\left|\int_\M S_1(\alpha)S_\lambda(\alpha)e(-n\alpha)\,d\alpha\right|
\ll X^{1-4\eta/9}L^{19/2}.
\end{equation}
Put $V(\beta)=\int_{X_1}^Xe(\beta u)\,du$.

\begin{fact}\label{fact:Euler-Maclaurin}
\cite[Lemma~2.8]{C9} Euler--Maclaurin summation gives
$|\mathcal T(\beta)-V(\beta)|\ll1+X|\beta|$.
\end{fact}

On $\M(q,a)$ this is $O(1+P/q)$. Therefore
\[
\int_\M\left|\frac{\mu(q)}{\varphi(q)}(\mathcal T-V)\right|^2
\ll\sum_{q\le P}\frac1{\varphi(q)}\cdot\frac{P}{qX}
\left(1+\frac Pq\right)^2
\ll\frac{P^3}{X}.
\]
Its mixed contribution with $S_\lambda$ is
\begin{equation}\label{eq:pole-discretization}
O(P^{3/2}L^{1/2}).
\end{equation}
Let $B_q=P/(qX)$ and $\Si(t)=\int_0^t(\sin u)/u\,du$. The remaining
pole term is
\begin{align}
R_{\mathrm{pole}}(n)
&=\sum_{q\le P}\frac{\mu(q)}{\varphi(q)}
\sum_{m\in I}\lambda(m)c_q(n-m)K_q(n-m),\label{eq:pole-term}\\
K_q(v)&=\frac1\pi\Bigl[\Si\bigl(2\pi B_q(X-v)\bigr)
+\Si\bigl(2\pi B_q(v-X_1)\bigr)\Bigr].
\label{eq:kernel}
\end{align}
This follows by integrating $V(\beta)e(-\beta v)$ over $[-B_q,B_q]$.
By \eqref{eq:prime-location}, both arguments have positive gaps at least
$X/8$. Here, since $\varepsilon_0>0$ is fixed and $X_1/X=X^{-\varepsilon_0}\to0$,
we have $X_1\le X/8$ for all sufficiently large $X$.
Differentiating \eqref{eq:kernel} gives
\begin{equation}\label{eq:kernel-derivative}
K_q'(v)=\frac1\pi\left(
\frac{\sin\bigl(2\pi B_q(v-X_1)\bigr)}{v-X_1}
-\frac{\sin\bigl(2\pi B_q(X-v)\bigr)}{X-v}
\right).
\end{equation}
On $[X/4,7X/8]$, each of $X-v$ and $v-X_1$ ranges within
$[X/8,7X/8]$. Hence
\[
\operatorname{Var}_{[X/4,7X/8]}K_q
=\int_{X/4}^{7X/8}|K_q'(v)|\,dv
\le\frac{2\log 7}{\pi}<2.
\]
The elementary estimates
$|\Si(t)-\pi/2|\le2/t$ and $\sup_{t\ge0}\Si(t)<1.86$ give
\begin{equation}\label{eq:kernel-bounds}
K_q(v)=1+O(q/P),\qquad \|K_q\|_\infty\le1.19,
\qquad \operatorname{Var}_{[X/4,7X/8]}K_q<2.
\end{equation}
We keep $K_q$ until the residue-class calculation is done.

\subsection{Reduced residue classes}

Only squarefree moduli contribute to \eqref{eq:pole-term}. All
non-squarefree arcs are already included in \eqref{eq:S1-removal}, since
there $\mu(q)=0$ and $S_1=S$. In the rest of the pole calculation, $q$ is
squarefree.

For $(r,q)=1$, partial summation and \eqref{eq:kernel-bounds} give
\begin{equation}\label{eq:kernel-distribution}
\begin{split}
\sum_{\substack{m\in I\\m\equiv r\pmod q}}\lambda(m)K_q(n-m)
={}&\frac1{\varphi(q)}
\sum_{\substack{m\in I\\(m,q)=1}}\lambda(m)K_q(n-m)\\
&+O(7\Delta_\lambda(q)).
\end{split}
\end{equation}
The bounded variation keeps a factor $P/q$ out of this estimate.
For squarefree $q$, multiplicativity gives
\begin{equation}\label{eq:Ramanujan-identities}
\begin{split}
\sum_{(r,q)=1}c_q(n-r)&=\mu(q)c_q(n),\\
\sum_{(r,q)=1}|c_q(n-r)|
&\le\sum_{r\bmod q}|c_q(n-r)|=2^{\omega(q)}\varphi(q).
\end{split}
\end{equation}
For a prime $p$, use $c_p(n)=p-1$ when $p\mid n$, and $c_p(n)=-1$
otherwise; the Chinese remainder theorem gives the product identities.
Consequently, the coprime part of \eqref{eq:pole-term} is
\begin{equation}\label{eq:coprime-pole}
\sum_{q\le P}\frac{\mu^2(q)c_q(n)}{\varphi(q)^2}
\sum_{\substack{m\in I\\(m,q)=1}}\lambda(m)K_q(n-m)
+O\left(\sum_{q\le Q_B}2^{\omega(q)}\Delta_\lambda(q)\right).
\end{equation}
Proposition~\ref{prop:BV9} bounds the error by $O(XL^{-B})$, since it
applies to this same weight, with both maxima, at level $Q_B\ge P$.

\subsection{Non-coprime classes}

The non-coprime part satisfies
\begin{equation}\label{eq:noncoprime-absolute}
|R_{\mathrm{nc}}(n)|
\le\sum_{q\le P}\frac{\mu^2(q)}{\varphi(q)}
\sum_{\substack{m\in I\\(m,q)>1}}
|\lambda(m)c_q(n-m)K_q(n-m)|.
\end{equation}
If a summand is nonzero, some prime $\ell\ge z$ divides both $m,q$.
Since $q<z^5$, there are at most four such primes in $q$.
Write $q=\ell s$. Squarefreeness gives $(\ell,s)=1$ and
\[
c_q(n-m)=c_\ell(n-m)c_s(n-m)=c_\ell(n)c_s(n-m).
\]

Also $|c_s(v)|\le(s,v)\le\sum_{d\mid s,\ d\mid v}d$.
For $d\mid s$ the conditions $\ell\mid m$ and $d\mid n-m$ are coprime,
and their count up to $X/4$ is at most $X/(\ell d)+1$.
Thus, termwise in the absolute sum,
\[
\begin{aligned}
\left|\sum_{\substack{m\in I\\\ell\mid m}}
\lambda(m)c_q(n-m)K_q(n-m)\right|
&\ll L|c_\ell(n)|\sum_{d\mid s}d
\left(\frac X{\ell d}+1\right)\\
&=L|c_\ell(n)|
\left(\frac{\tau(s)X}{\ell}+\sigma(s)\right).
\end{aligned}
\]
\[
\begin{aligned}
\text{RHS of \eqref{eq:noncoprime-absolute}}
&\ll L\sum_{z\le\ell\le P}\frac{|c_\ell(n)|}{\varphi(\ell)}
\sum_{s\le P/\ell}\frac{\tau(s)X/\ell+\sigma(s)}{\varphi(s)}\\
&\ll L^3\sum_{z\le\ell\le P}
\frac{|c_\ell(n)|}{\varphi(\ell)}\cdot\frac{X+P}{\ell}\\
&\ll \frac{XL^3}{z}=X^{9/10}L^3.
\end{aligned}
\]
Using $\sum_{s\le S}\tau(s)/\varphi(s)\ll L^2$ and
$\sigma(s)/\varphi(s)\ll L$, we obtain
\begin{align}
|R_{\mathrm{nc}}(n)|
&\ll L^3\sum_{z\le\ell\le P}
\frac{|c_\ell(n)|}{\varphi(\ell)}\frac{X+P}{\ell}
\ll XL^3/z=X^{9/10}L^3.
\label{eq:noncoprime-bound}
\end{align}
If $\ell\nmid n$, sum $1/(\ell\varphi(\ell))\ll1/z$.
If $\ell\mid n$, its contribution is at most $2X/z$, and at most ten
primes at least $z$ divide $n\le X=z^{10}$. This proves uniformity in
$n$.

We next remove $(m,q)=1$ from \eqref{eq:coprime-pole}. On the support of
$\lambda$, at most $4(X/z+1)$ non-coprime integers occur. For squarefree $q$, write $d=(q,n)$ and $q=dr$.
Then $|c_q(n)|=\varphi(d)$ and
\begin{equation}\label{eq:absolute-singular-series}
\sum_q\frac{\mu^2(q)|c_q(n)|}{\varphi(q)^2}
\le\sum_{d\mid n}\frac1{\varphi(d)}\sum_{r\ge1}\frac1{\varphi(r)^2}
\ll_\nu X^\nu.
\end{equation}
Together with $|\lambda K_q|\ll L$, this bounds the removal error by
\begin{equation}\label{eq:remove-coprimality}
O_\nu(X^{9/10+\nu})\qquad(\nu>0).
\end{equation}
The divisor-bound exponent is arbitrary, so it absorbs the logarithm.

\subsection{Completion of the singular series}

The same squarefree factorisation gives, for every $\nu>0$,
\begin{equation}\label{eq:singular-tails}
\sum_{q\le P}\frac{\mu^2(q)q|c_q(n)|}{\varphi(q)^2}\ll_\nu X^\nu,
\qquad
\sum_{q>P}\frac{\mu^2(q)|c_q(n)|}{\varphi(q)^2}\ll_\nu X^\nu/P.
\end{equation}
For the first bound, use
$\sum_{r\le Y}r/\varphi(r)^2\ll\log(2Y)$ and sum $d/\varphi(d)$ over $d\mid n$.
For the second, use $\sum_{r>Y}\varphi(r)^{-2}\ll1/Y$.
These elementary estimates follow by expanding the factors over squarefree
divisors. The resulting Euler products converge after the indicated
logarithmic factor is removed.

Both steps below use the absolute mass bound of Lemma~\ref{lem:weight-norms}, restated in \eqref{eq:mass-bound}:
\[
\sum_{m\in I}|\lambda(m)|\ll X,
\qquad |F_\lambda|\le\sum_{m\in I}|\lambda(m)|\ll X.
\]
By \eqref{eq:kernel-bounds}, replacing $K_q$ by $1$ costs
$O_\nu(X^{1+\nu}/P)$, since $\sum_m|\lambda(m)|\ll X$.
Completing the $q$-sum has the same cost, since $|F_\lambda|\ll X$.
The completed series is
\begin{equation}\label{eq:singular-series}
\begin{split}
\mathfrak S(n)
&=\sum_{q\ge1}\frac{\mu^2(q)c_q(n)}{\varphi(q)^2}\\
&=2\prod_{p>2}\left(1-\frac1{(p-1)^2}\right)
\prod_{\substack{p\mid n\\p>2}}\frac{p-1}{p-2}
\ge2C_2\ge1.32\qquad(2\mid n),
\end{split}
\end{equation}
where $C_2=\prod_{p>2}(1-(p-1)^{-2})$ and $2C_2=1.3203\ldots$.

\subsection{The major-arc asymptotic and positivity}

\begin{proposition}\label{prop:major-asymptotic}
Assume $H_{7/8}$ and the parameter choices of Section~\ref{sec:notation}.
Let $\lambda=(\log m)\rho(m)$ be the weight of Section~\ref{sec:weight}.
Thus it is supported on $z$-rough $m\in I$, with
$|\lambda|\le256L$, $\sum|\lambda|^2\ll XL$ and $\sum|\lambda|\ll X$.
Use Proposition~\ref{prop:BV9} for this weight at $Q_B\ge P$.
For $B>0$ and $0<\nu\le1/20$, uniformly for even $n\in(X/2,X]$,
\begin{equation}\label{eq:major-asymptotic}
\begin{split}
I_\M(n)=\mathfrak S(n)F_\lambda
+O_{\eta,B,\nu}\bigl(&X^{1-4\eta/9}L^{19/2}+P^{3/2}L^{1/2}+XL^{-B}\\
&+X^{9/10}L^3+X^{9/10+\nu}+X^{1+\nu}/P\bigr).
\end{split}
\end{equation}
In particular, $I_\M(n)=\mathfrak S(n)F_\lambda+o(X)$.
\end{proposition}

\begin{proof}
Use \eqref{eq:S1-removal} and \eqref{eq:pole-discretization}, and
evaluate the pole by \eqref{eq:coprime-pole} and
\eqref{eq:noncoprime-bound}. Then \eqref{eq:remove-coprimality} and
\eqref{eq:singular-tails} give exactly the six errors in
\eqref{eq:major-asymptotic}. At the stipulated $P$,
\begin{gather*}
P^{3/2}L^{1/2}\le X^{27/40-3\eta/2}L^{1/2},\qquad
X^{9/10+\nu}\le X^{19/20},\\
X^{1+\nu}/P\le2X^{11/20+\eta+\nu}\le2X^{61/100}.
\end{gather*}

Each error is $o(X)$, and this conclusion does not require a sign for
$F_\lambda$.
\end{proof}

\begin{corollary}\label{cor:major-positive}
Assume $H_{7/8}$. For all large $X$ and every even $n\in(X/2,X]$,
\[
I_\M(n)\ge0.025X.
\]
\end{corollary}

\begin{proof}
Proposition~\ref{prop:mass} gives $F_\lambda\ge0$ for large $X$.
By \eqref{eq:singular-series} and Proposition~\ref{prop:major-asymptotic},
\[
\mathfrak S(n)F_\lambda
\ge1.32F_\lambda
\ge1.32\Bigl(\frac{\ccert}{8}-o(1)\Bigr)X.
\]
The exact rational \eqref{eq:ccert} satisfies
$1.32\ccert/8=0.0250913982659\ldots>0.025$, and the positive margin
absorbs the error terms.
\end{proof}

\section{Proof of Theorem 2}\label{sec:proof-conditional}

\begin{proof}[Proof of Theorem~\ref{thm:conditional}]
Let $\varepsilon>0$ and use \eqref{eq:parameters}. At each large scale
$X$, choose $P$ as in \eqref{eq:P-window}. Corollary~\ref{cor:major-positive} supplies $c_1=0.025$ in
Proposition~\ref{prop:duality}.
Corollary~\ref{cor:minor-lambda} gives
$\mathcal S\ll X^{31/40+\eta/2}L^{9/2}$ at the same $P$, and both use the
same signed weight $\lambda$. Therefore
\[
|\mathcal E_X|
\ll\frac{\mathcal S^2L}{X}
\ll X^{31/20+\eta-1}L^{10}
=X^{11/20+\eta}L^{10}=X^{1-t}L^{10}.
\]
Choose $X_0=X_0(\varepsilon)$ so that this holds at every scale at least
$X_0$, and apply it to $Y=X/2^k$ with its own cutoff $P(Y)$.
The constants are uniform, and the geometric sum gives
\[
E(X)\le X_0+
\sum_{k\ge0}C_\varepsilon(X/2^k)^{11/20+\eta}(\log X)^{10}
\ll_\varepsilon X^{11/20+\eta}(\log X)^{10}
\ll_\varepsilon X^{11/20+\varepsilon}.
\]

Since $\eta<\varepsilon$, this is $O_\varepsilon(X^{11/20+\varepsilon})$.
\end{proof}

\section{Proof of Theorem 1}\label{sec:proof-main}

\begin{proof}[Proof of Theorem~\ref{thm:main}]
Theorem~\ref{thm:conditional}, together with the zero-free half-plane
$\Re s>7/8$ in \cite[Theorem~1.1]{OA}, gives the assertion.
\end{proof}

\appendix
\section{The exact rational mass certificate}\label{app:certificate}

This appendix gives the finite arithmetic for
\eqref{eq:cell-certificates}. The regions are exactly
\eqref{eq:loss-regions}, and the cell tests are \eqref{eq:U2-cell-tests}
and \eqref{eq:V4-cell-tests}. No limiting numerical integration is used.

\subsection{Integer form of the cell weights}

For a cell of dimension $k$ on the grid of size $n$, put
\[
r_i=2n+7j_i,\qquad r_i^+=r_i+7,\qquad
q_K=20n-\sum_{i=1}^k r_i^+,\qquad s_K=r_k^+.
\]
Thus $x_{i,0}=r_i/(20n)$ and $u_0(K)=q_K/s_K$. Write the rational
envelope value as $a_K/b_K$, using
\begin{equation}\label{eq:integer-envelope}
\frac{a_K}{b_K}=
\begin{cases}
1,&q_K<s_K,\\
s_K/q_K,&q_K\ge s_K\text{ and }125s_K\ge71q_K,\\
71/125,&q_K\ge s_K\text{ and }125s_K<71q_K.
\end{cases}
\end{equation}
Reduction to lowest terms is unnecessary, and \eqref{eq:cell-majorant}
is exactly
\begin{equation}\label{eq:integer-cell-weight}
W(K)=\frac{A_K}{B_K},\qquad
A_K=7^k(20n)a_K,\qquad
B_K=b_Kr_1\cdots r_{k-1}r_k^2.
\end{equation}
For $U_2$, take $(n,k)=(256,2)$ and $n>j_1\ge j_2\ge0$; the integer
tests are
\begin{equation}\label{eq:integer-U2-tests}
r_1<9n,\qquad r_1+r_2+14>11n,\qquad r_1+2r_2\le20n.
\end{equation}
For $V_4$, take $(n,k)=(64,4)$ and $n>j_1\ge j_2\ge j_3\ge j_4\ge0$;
the integer tests are
\begin{equation}\label{eq:integer-V4-tests}
r_1+r_2<9n,\quad r_1+r_2+r_3+21>11n,\quad
r_1+r_2+r_3+2r_4\le20n.
\end{equation}
These are the original rational cell tests after clearing denominators.

\subsection{Readable rational upper sums}

For each retained cell define the integer
\begin{equation}\label{eq:integer-ceiling}
b(K)=\left\lceil10^{18}W(K)\right\rceil
=\left\lfloor\frac{10^{18}A_K+B_K-1}{B_K}\right\rfloor.
\end{equation}
Every cell weight is rational and nonnegative, so
$W(K)\le b(K)/10^{18}$.
Tables~\ref{tab:U2-integers} and~\ref{tab:V4-integers} list sums of these
integers, with rows that partition the first-coordinate index $j_1$.
Each entry is obtained by the finite divisions and additions
in \eqref{eq:integer-envelope}--\eqref{eq:integer-ceiling}.

\begin{table}[htbp]
\caption{The $U_2$ upper sum. The final column has denominator $10^{18}$.}
\label{tab:U2-integers}
\centering\small
\renewcommand{\arraystretch}{1.2}
\begin{tabular}{@{}crr@{}}
\hline
$j_1$ & Retained cells & $\sum b(K)$\\
\hline
$0$--$127$ & $0$ & $0$\\
$128$--$159$ & $1056$ & $57250926781949766$\\
$160$--$191$ & $2773$ & $153887915136000733$\\
$192$--$223$ & $3376$ & $212249038575035418$\\
$224$--$255$ & $3888$ & $309273586283258073$\\
\hline
Total & $11093$ & $732661466776243990$\\
\hline
\end{tabular}
\end{table}

\begin{table}[htbp]
\caption{The $V_4$ upper sum. The final column has denominator $10^{18}$.}
\label{tab:V4-integers}
\centering\small
\renewcommand{\arraystretch}{1.2}
\begin{tabular}{@{}crr@{}}
\hline
$j_1$ & Retained cells & $\sum b(K)$\\
\hline
$0$--$7$ & $0$ & $0$\\
$8$--$15$ & $60$ & $307379357895441$\\
$16$--$23$ & $7919$ & $38542811400585796$\\
$24$--$31$ & $8263$ & $43932663439422582$\\
$32$--$39$ & $1926$ & $13933990209700055$\\
$40$--$47$ & $118$ & $1295654752060197$\\
$48$--$63$ & $0$ & $0$\\
\hline
Total & $18286$ & $98012499159664071$\\
\hline
\end{tabular}
\end{table}

The totals imply
\begin{equation}\label{eq:certificate-comparison}
\begin{aligned}
\mathbf U_2&\le\frac{732661466776243990}{10^{18}}
<\frac{732661466777}{10^{12}},\\
\mathbf V_4&\le\frac{98012499159664071}{10^{18}}
<\frac{98012508379}{10^{12}}.
\end{aligned}
\end{equation}
The positive differences after multiplication by $10^{18}$ are
$756010$ and $9219335929$. The final mass calculation uses the original
upper bounds on the right; the slightly smaller table bounds are not used.

\subsection{The four losses and the final subtraction}

Table~\ref{tab:mass-certificate} collects the rational coefficients of
$X/8$; the volume calculations for $L_4,U_4$ are proved in
Lemma~\ref{lem:volume-losses}.

\begin{table}[htbp]
\caption{Rational coefficients for the four discarded families.}
\label{tab:mass-certificate}
\centering\small
\renewcommand{\arraystretch}{2.45}
\begin{tabular}{@{}clc@{}}
\hline
Family & Calculation & Upper coefficient\\
\hline
$U_2$ & $256$-grid rectangle sum & $\dfrac{732661466777}{10^{12}}$\\
$L_4$ & $C_\omega10^5(1/20)^4/576$ & $\dfrac{71}{115200}$\\
$U_4$ & $C_\omega10^5(3/20)^4/1728$ & $\dfrac{213}{12800}$\\
$V_4$ & $64$-grid rectangle sum & $\dfrac{98012508379}{10^{12}}$\\
\hline
\end{tabular}
\end{table}

Exact subtraction gives
\[
1-\frac{732661466777}{10^{12}}-\frac{71}{115200}
-\frac{213}{12800}-\frac{98012508379}{10^{12}}
=\frac{342155430899}{2250000000000}.
\]
In particular,
\[
\ccert-\frac3{20}
=\frac{4655430899}{2250000000000}>0.
\]
The envelope inequality is also certified rationally:
\[
\frac{71}{125}\sum_{j=0}^4\frac{(71/125)^j}{j!}
=\frac{733930882351}{732421875000}>1.
\]
The accompanying script \texttt{cert1120.py} evaluates the original cell
sums with exact rational arithmetic. Its output confirms the cell counts,
the two upper bounds, the two volume coefficients and the displayed value
of $\ccert$. The analytic reduction to these finite sums is proved in
Section~\ref{sec:mass}. The script checks only the arithmetic and supplies
no asymptotic estimate.

\section*{Statements and declarations}

\textbf{Competing interests.} The author has no competing interests.

\textbf{Funding.} No funding was received for this work.

\textbf{Data availability.} The rational certificate is specified in
Appendix~\ref{app:certificate}. The accompanying script
\texttt{cert1120.py} and its output reproduce the finite sums.

\begin{thebibliography}{10}

\bibitem{Da}
H.~Davenport, \emph{Multiplicative Number Theory}, 2nd ed.,
revised by H.~L.~Montgomery, Graduate Texts in Mathematics \textbf{74},
Springer, New York, 1980.

\bibitem{C4}
P.~X.~Gallagher, Invent. Math. \textbf{11} (1970), 329--339, Lemma~1.

\bibitem{HB}
D.~R.~Heath-Brown, Canad. J. Math. \textbf{31} (1979), 231--240,
Theorem~1 and Theorem~2(3).

\bibitem{IK}
H.~Iwaniec and E.~Kowalski, \emph{Analytic Number Theory},
Section~3.4 and Theorem~7.13.

\bibitem{C2c}
H.~L.~Montgomery, \emph{Topics in Multiplicative Number Theory},
Lecture Notes in Mathematics \textbf{227}, Theorem~12.2.

\bibitem{OA}
OpenAI, \emph{The Quasi-Riemann Hypothesis: A Zero-Free Half-Plane
$\Re s>7/8$}, OpenAI Math Release preprint, 30 September 2026,
family 003, Theorem~1.1.
Primary source PDF: \url{https://github.com/openai/math/blob/adc7f1241b42e322a6451854ab7e4b4c146bf78a/preprints/The-Quasi-Riemann-Hypothesis-September-30-2026/paper.pdf},\newline
commit \texttt{adc7f1241b42e322a6451854ab7e4b4c146bf78a};\newline
PDF SHA-256 \texttt{8fe93046f8cf5ef1ba5969c89addc02d76311adc4ee907509ff9cd96f7ec99e7}.

\bibitem{Pi}
J.~Pintz, \emph{A new explicit formula in the additive theory of primes
with applications I. The explicit formula for the Goldbach problem and the
Generalized Twin Prime Problem}, Acta Arith. \textbf{210} (2023), 53--94,
doi:10.4064/aa220728-31-3.
Equation and theorem numbers cited here are from the preprint version dated
16 April 2018.

\bibitem{PA}
J.~Pintz and I.~Alth\"ofer, \emph{The Goldbach exceptional set: an exponent
$7/12$ and conditional refinement}, preprint, arXiv:2610.15559, version~1
(11 October 2026),
Theorem~1 and Corollary~2.

\bibitem{C9}
R.~C.~Vaughan, \emph{The Hardy--Littlewood Method}, 2nd ed.,
Lemmas~2.2 and~2.8.

\end{thebibliography}

\end{document}
